%%%%%%%%%%%%%%%%%%%%%%%%%%%%Nv-Sujet%%%%%%%%%%%%%%%%%%%%%%%%%%%%%%%%%%%%%%%%%%

\sujetn{%
	nomauteur=Michaux,
	prenomauteur=Henri,
	initialeprenom =H,
	titre = {Qui je fus},
	semestre = \HQ,
	lieu =Sujet 0,
	annee =2020,
	type ={Non-proposé},
	jour =0,
	datepub =1927,
	ajoutref =,
	genre =\poesie,
	corrige =Oui,
	liencorrige={sujet0_michaux},
	typequn = \litt,
	typeqdeux = \phil,
	qun = {Comment ce poème dit-il la violence, et qu'en dit-il~?},
	qdeux ={La violence échappe-t-elle à notre compréhension~?},
	themeprogrammequn=HV,
	themeprogrammeqdeux=HV}
{\setlength{\parskip}{0pt}
{\begin{center} LE GRAND COMBAT\end{center}}

\bigskip

Il l'emparouille et l'endosque contre terre~;

Il le rague et le roupète jusqu'à son drâle~;

Il le pratèle et le libucque et lui barufle les ouallais~;

Il le tocarde et le marmine,

Le manage rape à ri et ripe à ra.

Enfin il l'écorcobalisse.

L'autre hésite, s'espudrine, se défaisse, se torse et se ruine.

C'en sera bientôt fini de lui~;

Il se reprise et s'emmargine… mais en vain

Le cerceau tombe qui a tant roulé.

Abrah~! Abrah~! Abrah~!

Le pied a failli~!

Le bras a cassé~!

Le sang a coulé~!

Fouille, fouille, fouille,

Dans la marmite de son ventre est un grand secret

Mégères alentour qui pleurez dans vos mouchoirs~;

On s'étonne, on s'étonne, on s'étonne

Et vous regarde

On cherche aussi, nous autres, le Grand Secret.

\setlength{\parskip}{6pt}
\theme{combat}\theme{néologisme}\theme{hermétisme}\theme{tristesse}\theme{souffrance}}

%%%%%%%%%%%%%%%%%%%%%%%%%%%%Nv-Sujet%%%%%%%%%%%%%%%%%%%%%%%%%%%%%%%%%%%%%%%%%%

\sujetn{%
	nomauteur=Arendt,
	prenomauteur=Hannah,
	initialeprenom =H,
	particule=,
	titre = {Condition de l'homme moderne},
	semestre = \HQ,
	lieu =Sujet 0,
	annee =2020,
	type =Non-proposé,
	jour =0,
	datepub =1958,
	ajoutref ={traduit de l'anglais par Georges \textsc{Fradier}},
	genre =\essaip,
	corrige =Oui,
	liencorrige={sujet0_arendt},
	typequn = \phil,
	typeqdeux = \litt,
	qun = {Par une lecture attentive du texte et de son argumentation, expliquez pourquoi la question de l'~homme futur~» n'est pas une question purement technique, mais bien une question de nature politique.},
	qdeux ={«~C'est une question politique primordiale que l'on ne peut guère, par conséquent, abandonner aux professionnels de la science ni à ceux de la politique~». Que peuvent apporter à la réflexion sur cette question les arts et la littérature~?},
	themeprogrammequn=HL,
	themeprogrammeqdeux=HVL}
{La nature terrestre, pour autant que l'on sache, pourrait bien être la
seule de l'univers à procurer aux humains un habitat où ils puissent
se mouvoir et respirer sans effort et sans artifice. L'artifice
humain du monde sépare l'existence humaine de tout milieu purement
animal, mais la vie elle-même est en dehors de ce monde artificiel,
et par la vie l'homme demeure lié à tous les autres organismes
vivants. Depuis quelque temps, un grand nombre de recherches
scientifiques s'efforcent de rendre la vie artificielle elle aussi,
et de couper le dernier lien qui maintient encore l'homme parmi les
enfants de la nature. C'est le même désir d'échapper à
l'emprisonnement terrestre qui se manifeste dans les essais de
création en éprouvette, dans le vœu de combiner «~au microscope le
plasma germinal provenant de personnes aux qualités garanties, afin
de produire des êtres supérieurs~» et «~de modifier (leurs) tailles,
formes et fonctions~»\footnote{Hannah Arendt fait ici référence à des formules qui ont été utilisées dans l'espace public et dans les médias de l'époque, lors du lancement par l'Union soviétique du premier satellite artificiel, Spoutnik 1, le 4 octobre 1957.}~; et je soupçonne que l'envie d'échapper à
la condition humaine expliquerait l'espoir de prolonger la durée de
l'existence fort au-delà de cent ans, limite jusqu'ici admise.

Cet homme futur, que les savants produiront, nous disent-ils, en un
siècle pas davantage, paraît en proie à la révolte contre l'existence
humaine telle qu'elle est donnée, cadeau venu de nulle part
(laïquement parlant) et qu'il veut pour ainsi dire échanger contre un
ouvrage de ses propres mains. Il n'y a pas de raison de douter que
nous soyons capables de faire cet échange, de même qu'il n'y a pas de
raison de douter que nous soyons capables à présent de détruire toute
vie organique sur terre. La seule question est de savoir si nous
souhaitons employer dans ce sens nos nouvelles connaissances
scientifiques et techniques, et l'on ne saurait en décider par des
méthodes scientifiques. C'est une question politique primordiale que
l'on ne peut guère, par conséquent, abandonner aux professionnels de
la science ni à ceux de la politique.
\theme{vie (et mort)}\theme{naturel (et artificiel)}\theme{eugénisme}\theme{biogénétique}\theme{transhumanisme}\theme{technique!maîtrise technique de la nature}\theme{technique!et sciences}\theme{immortalité}\theme{vie (et mort)!sens de la vie}\theme{nature!nature humaine}}

%%%%%%%%%%%%%%%%%%%%%%%%%%%%Nv-Sujet%%%%%%%%%%%%%%%%%%%%%%%%%%%%%%%%%%%%%%%%%%

\sujetn{%
	nomauteur=Veil,
	prenomauteur=Simone,
	initialeprenom =S,
	particule=,
	titre = {Une vie},
	semestre = \HQ,
	lieu ={Amérique du Nord},
	annee =2021,
	type =,
	jour =1,
	datepub =2007,
	ajoutref =,
	genre =\autobiographie,
	corrige =Oui,
	liencorrige={amnord_2021_1_2},
	typequn = \litt,
	typeqdeux = \phil,
	qun = {De quelles manières Simone Veil fait-elle de la mémoire des vivants
aussi une mémoire des morts~?},
	qdeux ={Témoigner de la violence, est-ce un besoin ou un devoir~?},
	themeprogrammequn=HV,
	themeprogrammeqdeux=HV}
{Parler de la Shoah, et comment~; ou bien ne pas en parler, et pourquoi
? Éternelle question. Le romancier israélien Aharon Appelfeld a écrit
plusieurs livres superbes, notamment \textit{Histoire d'une vie}, où
il raconte son évasion du camp, alors qu'il a dix ans, et ses trois
ans de cache dans la forêt ukrainienne. Il vient de publier trois
discours prononcés en Israël. C'est un livre bouleversant dans lequel
il analyse la Shoah en expliquant que ceux qui en ont été les
victimes ne s'en sortent jamais. À sa lecture, je me suis rendu
compte qu'au fond nous aurons toujours vécu avec cela. Certains
répugnent à l'évoquer. D'autres ont besoin d'en parler. Mais tous
vivent avec.

Appelfeld énonce les raisons pour lesquelles on ne peut plus s'en
détacher. Elles sont terribles, et marquent la différence de nature
avec la situation des résistants. Eux sont dans la position des
héros, leur combat les couvre d'une gloire qu'accroît encore
l'emprisonnement dont ils l'ont payée~; ils avaient choisi leur
destin. Mais nous, nous n'avions rien choisi. Nous n'étions que des
victimes honteuses, des animaux tatoués. Il nous faut donc vivre avec
ça, et que les autres l'acceptent. Tout ce qu'on peut dire, écrire,
filmer sur l'Holocauste n'exorcise rien. La Shoah est omniprésente.
Rien ne s'efface~; les convois, le travail, l'enfermement, les
baraques, la maladie, le froid, le manque de sommeil, la faim, les
humiliations, l'avilissement, les coups, les cris… non, rien ne peut
ni ne doit être oublié. Mais au-delà de ces horreurs, seuls importent
les morts. La chambre à gaz pour les enfants, les femmes, les
vieillards, pour ceux qui attrapent la gale, qui clopinent, qui ont
mauvaise mine~; et pour les autres, la mort lente. Deux mille cinq
cents survivants sur soixante-dix-huit mille Juifs français déportés.
Il n'y a que la Shoah. L'atmosphère de crématoire, de fumée et de
puanteur de Birkenau, je ne l'oublierai jamais. Là-bas, dans les
plaines allemandes et polonaises, s'étendent désormais des espaces
dénudés sur lesquels règne le silence~; c'est le poids effrayant du
vide que l'oubli n'a pas le droit de combler, et que la mémoire des
vivants habitera toujours.
\theme{Shoah}\theme{camp de concentration}\theme{victime}\theme{mémoire}\theme{oubli}\theme{déportation}}

%%%%%%%%%%%%%%%%%%%%%%%%%%%%Nv-Sujet%%%%%%%%%%%%%%%%%%%%%%%%%%%%%%%%%%%%%%%%%%

\sujetn{%
	nomauteur=Jonas,
	prenomauteur=Hans,
	initialeprenom =H,
	particule=,
	titre = {Une éthique pour la nature},
	semestre = \HQ,
	lieu ={Amérique du Nord},
	annee =2021,
	type =,
	jour =2,
	datepub =1987,
	ajoutref ={«~Technique, liberté, obligation~», trad.~Sylvie \textsc{Courtine-Denamy}, Paris, Arthaud},
	genre =\essaip,
	corrige =Oui,
	liencorrige={amnord_2021_2_2},
	typequn = \phil,
	typeqdeux = \litt,
	qun = {Pourquoi, d'après Hans Jonas, l'homme est-il devenu un danger pour lui-même~?},
	qdeux ={La littérature nous rend-elle sensible à l'importance de la nature pour l'homme~?},
	themeprogrammequn=HL,
	themeprogrammeqdeux=EXSHL}
{\label{jonas-amérique2021}
Le danger qui nous menace actuellement vient-il encore du dehors~?
Provient-il de l'élément sauvage que nous devons maîtriser grâce aux
formations artificielles de la culture~? C'est encore parfois le cas,
mais un flot nouveau et plus dangereux se déchaîne maintenant de
l'intérieur même et se précipite, détruisant tout sur son passage, y
compris la force débordante de nos actions qui relèvent de la
culture. C'est désormais à partir de nous que s'ouvrent les trouées
et les brèches à travers lesquelles notre poison se répand sur le
globe terrestre, transformant la nature tout entière en un cloaque\footnote{Cloaque~: lieu très sale et malsain prévu pour recevoir les détritus de toutes sortes.}
pour l'homme. Ainsi les fronts se sont-ils inversés. Nous devons
davantage protéger l'océan contre nos actions que nous protéger de
l'océan. Nous sommes devenus un plus grand danger pour la nature que
celle-ci ne l'était autrefois pour nous. Nous sommes devenus
extrêmement dangereux pour nous-mêmes et ce, grâce aux réalisations
les plus dignes d'admiration que nous avons accomplies pour assurer
la domination de l'homme sur les choses. C'est nous qui constituons
le danger dont nous sommes actuellement cernés et contre lequel nous
devons désormais lutter. Il s'agit là de quelque chose de
radicalement nouveau~: aucune des obligations que nous connaissons
n'est jamais née d'une impulsion salvatrice commune.
\theme{naturel (et artificiel)}\theme{culture!comme maîtrise de la nature}\theme{nature!environnement}\theme{écologie}\theme{technique!maîtrise technique de la nature}}

%%%%%%%%%%%%%%%%%%%%%%%%%%%%Nv-Sujet%%%%%%%%%%%%%%%%%%%%%%%%%%%%%%%%%%%%%%%%%%

\sujetn{%
	nomauteur=Ferrat,
	prenomauteur=Jean,
	initialeprenom =J,
	particule=,
	titre = {Album Nuit et Brouillard},
	semestre = \HQ,
	lieu =Asie,
	annee =2021,
	type =,
	jour =2,
	datepub =1963,
	ajoutref = Barclay,
	genre =\chanson,
	corrige =,
	liencorrige={asi_2021_2_2},
	typequn = \litt,
	typeqdeux = \phil,
	qun = {Comment ce texte vient-il restaurer l'humanité mise en péril des déportés~?},
	qdeux ={Comment résiste-t-on à la déshumanisation~?},
	themeprogrammequn=HV,
	themeprogrammeqdeux=HVL}
{\chapeau{
Jean Ferrat (1930-2010) compose cette chanson, \textup{Nuit et
brouillard}, en 1963. «~Nuit et brouillard~» est le nom d'un décret
des autorités nazies, signé le 7 décembre 1941 et condamnant les
opposants à la déportation en Allemagne et à la mort.}

\setlength{\parskip}{0pt}

\begin{center}
Nuit et brouillard \end{center}


Ils étaient vingt et cent, ils étaient des milliers

Nus et maigres, tremblants, dans ces wagons plombés

Qui déchiraient la nuit de leurs ongles battants

Ils étaient des milliers, ils étaient vingt et cent

\bigskip

Ils se croyaient des hommes, n'étaient plus que des nombres

Depuis longtemps leurs dés avaient été jetés

Dès que la main retombe il ne reste qu'une ombre

Ils ne devaient jamais plus revoir un été

\bigskip

La fuite monotone et sans hâte du temps

Survivre encore un jour, une heure, obstinément

Combien de tours de roues, d'arrêts et de départs

Qui n'en finissent pas de distiller l'espoir

\bigskip

Ils s'appelaient Jean-Pierre, Natacha ou Samuel

Certains priaient Jésus, Jéhovah ou Vichnou

D'autres ne priaient pas, mais qu'importe le ciel

Ils voulaient simplement ne plus vivre à genoux

\bigskip

Ils n'arrivaient pas tous à la fin du voyage

Ceux qui sont revenus peuvent-ils être heureux

Ils essaient d'oublier, étonnés qu'à leur âge

Les veines de leurs bras soient devenues si bleues

\bigskip

Les Allemands guettaient du haut des miradors

La lune se taisait comme vous vous taisiez

En regardant au loin, en regardant dehors

Votre chair était tendre à leurs chiens policiers

\bigskip

On me dit à présent que ces mots n'ont plus cours

Qu'il vaut mieux ne chanter que des chansons d'amour

Que le sang sèche vite en entrant dans l'histoire

Et qu'il ne sert à rien de prendre une guitare

\bigskip

Mais qui donc est de taille à pouvoir m'arrêter~?
\nopagebreak\par
L'ombre s'est faite humaine, aujourd'hui c'est l'été
\nopagebreak\par
Je twisterais\footnote{«~twister~»~: tordre, tortiller, retourner (le
twist est une danse populaire des années 1960, qui se danse sur une
musique de rock and roll, par une rotation des jambes et du bassin).}
les mots s'il fallait les twister
\nopagebreak\par
Pour qu'un jour les enfants sachent qui vous étiez

\bigskip

Vous étiez vingt et cent, vous étiez des milliers

Nus et maigres, tremblants, dans ces wagons plombés

Qui déchiriez la nuit de vos ongles battants

Vous étiez des milliers, vous étiez vingt et cent
\setlength{\parskip}{6pt}
\theme{déportation}\theme{Shoah}\theme{déshumanisation}\theme{vie (et mort)}\theme{temps}\theme{histoire}\theme{mémoire}}

%%%%%%%%%%%%%%%%%%%%%%%%%%%%Nv-Sujet%%%%%%%%%%%%%%%%%%%%%%%%%%%%%%%%%%%%%%%%%%

\sujetn{%
	nomauteur=Weil,
	prenomauteur=Simone,
	initialeprenom =S,
	particule=,
	titre = {L'Iliade ou le poème de la force},
	semestre = \HQ,
	lieu =Centres étrangers,
	annee =2021,
	type =,
	jour =1,
	datepub =1941,
	ajoutref =,
	genre =\essaip,
	corrige =,
	liencorrige={cea_2021_1_2},
	typequn = \phil,
	typeqdeux = \litt,
	qun = {D'après ce texte, pourquoi la violence de la guerre nous
déshumanise-t-elle~?},
	qdeux ={Que peut dire la littérature confrontée à une violence sans issue~?},
	themeprogrammequn=HVL,
	themeprogrammeqdeux=HV}
{Que des hommes aient pour avenir la mort, cela est contre nature. Dès
que la pratique de la guerre a rendu sensible la possibilité de mort
qu'enferme chaque minute, la pensée devient incapable de passer d'un
jour à son lendemain sans traverser l'image de la mort. L'esprit est
alors tendu comme il ne peut souffrir de l'être que peu de temps~;
mais chaque aube nouvelle amène la même nécessité~; les jours ajoutés
aux jours font des années. L'âme souffre violence tous les jours.
Chaque matin l'âme se mutile de toute aspiration, parce que la pensée
ne peut pas voyager dans le temps sans passer par la mort. Ainsi la
guerre efface toute idée de but, même l'idée des buts de la guerre.
Elle efface la pensée même de mettre fin à la guerre. La possibilité
d'une situation si violente est inconcevable tant qu'on n'y est pas~;
la fin en est inconcevable quand on y est. Ainsi l'on ne fait rien
pour amener cette fin. Les bras ne peuvent pas cesser de tenir et de
manier les armes en présence d'un ennemi armé~; l'esprit devrait
combiner pour trouver une issue~; il a perdu toute capacité de rien
combiner à cet effet. Il est occupé tout entier à se faire violence.
Toujours parmi les hommes, qu'il s'agisse de servitude ou de guerre,
les malheurs intolérables durent par leur propre poids et semblent
ainsi du dehors faciles à porter~; ils durent parce qu'ils ôtent les
ressources nécessaires pour en sortir. Néanmoins l'âme soumise à la
guerre crie vers la délivrance~; mais la délivrance même lui apparaît
sous une forme tragique, extrême, sous la forme de la destruction.
Une fin modérée, raisonnable, laisserait à nu pour la pensée un
malheur si violent qu'il ne peut être soutenu même comme souvenir. La
terreur, la douleur, l'épuisement, les massacres, les compagnons
détruits, on ne croit pas que toutes ces choses puissent cesser de
mordre l'âme si l'ivresse de la force n'est venue les noyer.
\theme{guerre}\theme{temps}\theme{pensée!anéantissement de la pensée}\theme{souffrance}\theme{déshumanisation}\theme{armes}}

%%%%%%%%%%%%%%%%%%%%%%%%%%%%Nv-Sujet%%%%%%%%%%%%%%%%%%%%%%%%%%%%%%%%%%%%%%%%%%

\sujetn{%
	nomauteur=Guilloux,
	prenomauteur=Louis,
	initialeprenom =L,
	particule=,
	titre = {Le Sang noir},
	semestre = \HQ,
	lieu = Centres étrangers,
	annee =2021,
	type =,
	jour =2,
	datepub =1935,
	ajoutref =,
	genre =\romanl,
	corrige =,
	liencorrige={cea_2021_2_1},
	typequn = \litt,
	typeqdeux = \phil,
	qun = {Comment le personnage prend-il conscience de la violence du monde et comment le texte en rend-il compte~?},
	qdeux ={Pourquoi est-il dangereux de nier l'existence de la violence dans
l'Histoire~?},
	themeprogrammequn=HV,
	themeprogrammeqdeux=HV}
{\chapeau{
Pendant la Première Guerre mondiale, M. Marchandeau vient d'apprendre
que son fils Pierre, soldat ayant refusé d'obéir aux ordres,
s'apprête à être fusillé pour cette raison.}

Comme tout le monde, M. Marchandeau avait souvent compulsé, d'une main
parfois distraite, ces
\textit{illustrés}\footnote{«~illustré~»~:
magazine comportant des images.}
 de la guerre qui offraient au monde un tel résumé d'horreurs qu'il ne
semblait pas croyable que personne en pût supporter la vue. Dans ces
illustrés, dont certains se vantaient de payer n'importe quel prix
les documents intéressants, il lui était arrivé de tomber sur les
images d'une exécution capitale~: espion passé par les
armes\footnote{«~passé par les armes~»~: exécuté.}. L'homme, la tête
basse, les mains liées, une dernière cigarette aux lèvres, marchait
entouré de ses bourreaux, et M. Marchandeau avait remarqué qu'il s'en
trouvait toujours un pour sourire. C'était à croire qu'il ne pouvait
y avoir d'exécution capitale sans ce sourire-là~! Qui donc tout à
l'heure sourirait~?

Venait ensuite l'exécution proprement dite~: l'homme, à genoux devant
le poteau, les yeux bandés. Ensuite enfin, et pour conclure, le
défilé des troupes devant le cadavre.

Il avait regardé ces images non sans émotion, mais avec le sentiment
que cela ne le concernait pas directement, que ces choses atroces se
passaient dans un univers sans rapport avec le sien, si paisible, que
bien sûrement il ne serait jamais fusillé, lui ni personne qu'il
connût. Or…

Il lui arrivait, comme à tant d'autres, une aventure à laquelle il
n'était pas préparé~: il était au spectacle, commodément installé
dans un fauteuil, et voilà qu'on le priait durement de vider son
siège, de grimper en scène, d'y traîner avec lui sa femme et son
fils. Il n'avait pas prévu cela. Naïvement, jusqu'au 2 août 1914
\footnote{«~2 août 1914~»~: le jour de la mobilisation générale, au
cours duquel les soldats sont appelés au front.}, il avait pris la
vie pour un conte. On exigeait aujourd'hui, fouet en main, qu'il prît
au jeu une part active, sans même lui demander s'il avait au moins
appris un petit bout de rôle, s'il savait en quoi consistait le
scénario dans son ensemble et au bénéfice de qui était monté ce
gala\footnote{«~gala~»~: grande fête officielle.}~? Mais il ne savait
rien. Il voyait seulement qu'il ne s'agissait plus de spectacle du
tout, que la comédie tournait au drame – au vrai drame – que la balle
était une vraie balle, l'épée vraiment teintée de sang, le mort un
vrai mort.

On fusillait les espions~: soit~! Mais on ne lui avait pas dit qu'on
fusillait aussi les insurgés\footnote{«~insurgés~»~: ceux qui se
révoltent.}, ni même qu'il y en eût. On lui avait fait croire que
tout allait «~à merveille~» et que ces milliers de jeunes gens jetés
au fumier acceptaient joyeusement leur mort. Il s'était laissé duper
sans penser une seconde que la machine meurtrière pouvait aussi se
retourner contre lui et contre son fils. Il avait laissé faire, il
avait consenti. Il était complice, hélas~! de ce sourire qui tout à
l'heure accompagnerait Pierre au poteau, complice des prières qu'un
tendre aumônier\footnote{«~aumônier~»~: prêtre chargé de recueillir
les dernières volontés d'un mourant.} ne manquerait pas de prodiguer
à son fils afin que tout soit en règle et la mort bien parée.
\theme{guerre!Première Guerre mondiale}\theme{désobéissance}\theme{exécution}\theme{armée}\theme{parent}\theme{vie (et mort)}\theme{esthétique!de la violence}\theme{armes}}

%%%%%%%%%%%%%%%%%%%%%%%%%%%%Nv-Sujet%%%%%%%%%%%%%%%%%%%%%%%%%%%%%%%%%%%%%%%%%%

\sujetn{%
	nomauteur=Antelme,
	prenomauteur=Robert,
	initialeprenom =R,
	particule=,
	titre = {L'espèce humaine},
	semestre = \HQ,
	lieu =Métropole,
	annee =2021,
	type =Candidats libres,
	jour =1,
	datepub =1947,
	ajoutref =,
	genre =\romanl,
	corrige =,
	liencorrige={metcal_2021_1_2},
	typequn = \litt,
	typeqdeux = \phil,
	qun = {Qu'apporte à la réflexion de Robert Antelme sa méditation poétique sur la nature~?},
	qdeux ={L'espèce humaine a-t-elle besoin de faire l'expérience de la violence
pour éprouver son unité~?},
	themeprogrammequn=EXSHV,
	themeprogrammeqdeux=HV}
{\chapeau{
Publiée en 1947, \emph{L'Espèce humaine} est un récit de
l'expérience des camps de concentration et d'extermination pendant la
Deuxième Guerre mondiale. Résistant, Robert Antelme a été arrêté et
déporté d'abord à Buchenwald puis dans un camp de travail à
Gandersheim. L'extrait proposé se situe au moment où prisonniers et
officiers SS fuient sur les routes d'Allemagne l'avancée des armées
alliées. Après une longue journée de marche, les prisonniers se
retrouvent dans «~une grande maison isolée~» pour passer la nuit,
tous entassés sur le plancher.}

Dehors, la vallée est noire. Aucun bruit n'en arrive. Les chiens
dorment d'un sommeil sain et repu. Les arbres respirent calmement.
Les insectes nocturnes se nourrissent dans les prés. Les feuilles
transpirent, et l'air se gorge d'eau. Les prés se couvrent de rosée
et brilleront tout à l'heure au soleil. Ils sont là, tout près, on
doit pouvoir les toucher, caresser cet immense pelage. Qu'est-ce qui
se caresse et comment caresse-t-on~? Qu'est-ce qui est doux aux
doigts, qu'est-ce qui est seulement à être caressé~?

Jamais on n'aura été aussi sensible à la santé de la nature. Jamais on
n'aura été aussi près de confondre avec la toute-puissance l'arbre
qui sera sûrement encore vivant demain. On a oublié tout ce qui meurt
et qui pourrit dans cette nuit forte, et les bêtes malades et seules.
La mort a été chassée par nous des choses de la nature, parce que
l'on n'y voit aucun génie qui s'exerce contre elles et les poursuive.
Nous nous sentons comme ayant pompé tout pourrissement possible. Ce
qui est dans cette salle apparaît comme la maladie extraordinaire, et
notre mort ici comme la seule véritable. Si ressemblants aux bêtes,
toute bête nous est devenue somptueuse~; si semblables à toute plante
pourrissante, le destin de cette plante nous paraît aussi luxueux que
celui qui s'achève par la mort dans le lit. Nous sommes au point de
ressembler à tout ce qui ne se bat que pour manger et meurt de ne pas
manger, au point de nous niveler sur une autre espèce, qui ne sera
jamais nôtre et vers laquelle on tend~; mais celle-ci qui vit du
moins selon sa loi authentique – les bêtes ne peuvent pas devenir
plus bêtes – apparaît aussi somptueuse que la nôtre «~véritable~»
dont la loi peut être aussi de nous conduire ici. Mais il n'y a pas
d'ambiguïté, nous restons des hommes, nous ne finirons qu'en hommes.
La distance qui nous sépare d'une autre espèce reste intacte, elle
n'est pas historique. C'est un rêve SS de croire que nous avons pour
mission historique de changer d'espèce, et comme cette mutation se
fait trop lentement, ils tuent. Non, cette maladie extraordinaire
n'est autre chose qu'un moment culminant de l'histoire des hommes. Et
cela peut signifier deux choses~: d'abord que l'on fait l'épreuve de
la solidité de cette espèce, de sa fixité. Ensuite, que la variété
des rapports entre les hommes, leur couleur, leurs coutumes, leur
formation en classes masquent une vérité qui apparaît ici éclatante,
au bord de la nature, à l'approche de nos limites~: il n'y a pas des
espèces humaines, il y a une espèce humaine. C'est parce que nous
sommes des hommes comme eux que les SS seront en définitive
impuissants devant nous. C'est parce qu'ils auront tenté de mettre en
cause l'unité de l'espèce qu'ils seront finalement écrasés.
\theme{camp de concentration}\theme{nature!environnement}\theme{vie (et mort)}\theme{humain!et animal}\theme{déshumanisation}\theme{eugénisme}\theme{nature!nature humaine}\theme{déportation}\theme{Shoah}\theme{SS}}


%%%%%%%%%%%%%%%%%%%%%%%%%%%%Nv-Sujet%%%%%%%%%%%%%%%%%%%%%%%%%%%%%%%%%%%%%%%%%%

\sujetn{%
	nomauteur=Hyvernaud,
	prenomauteur=Georges,
	initialeprenom =G,
	particule=,
	titre = {La peau et les os},
	semestre = \HQ,
	lieu =Métropole,
	annee =2021,
	type =Candidats libres,
	jour =2,
	datepub =1949,
	ajoutref =,
	genre =\romanl,
	corrige =,
	liencorrige={metcal_2021_2_2},
	typequn = \litt,
	typeqdeux = \phil,
	qun = {À quelles formes diverses de violence le narrateur est-il exposé~?},
	qdeux ={L'expérience de la souffrance est-elle incommunicable~?},
	themeprogrammequn=HV,
	themeprogrammeqdeux=EXSHV}
{\chapeau{Prisonnier de guerre en Allemagne pendant la Seconde Guerre
mondiale, Georges Hyvernaud raconte ici son retour à la vie
familiale, après sa détention dans un camp de travail.}

Après que chacun a bien parlé de soi, la famille se rappelle pourtant
ma présence. Vous autres aussi, dans vos camps, vous en baviez, dit
la Famille. Forcément, on en bavait. Les têtes se tournent vers moi,
c'est mon tour. La Famille veut savoir ce que nous mangions, si les
gardiens nous maltraitaient. Raconte un peu, demande Louise, le type
qui s'est évadé dans une poubelle. Oh oui, raconte, implore la
Famille. Je me fais l'effet d'être encore le petit garçon à qui on
imposait de réciter au dessert \textit{La Mendiante}, d'Eugène
Manuel\footnote{Eugène Manuel~: poète et homme politique français
(1823-1901), dont l'œuvre est influencée par le romantisme, la poésie
parnassienne et le naturalisme.}. Je me résigne~: Eh bien, voilà,
c'est un type qui…

Mes souvenirs, dans ces moments où je suis bien encastré dans la paix
compacte de la Famille, c'est curieux comme ils perdent de leur
mordant et de leur autorité. Ils sont sans force, ils n'ont même plus
l'air vrai. Pas moyen de croire à ça quand on regarde Ginette servir
le café en prenant soin de ne pas tacher la nappe. Quand on regarde
Merlandon, le Vétérinaire, l'Oncle. Existences indiscutables et
invincibles comme celle des choses. Comme celle du petit berger de
bronze sur son napperon de dentelle – la même dignité, la même
puissance sourde. Cette solidité repousse et nie les souvenirs. Au
contact de la réalité des dimanches familiaux, l'humiliation et le
désespoir ne font plus qu'un jeu d'ombres improbables, une espèce de
cinéma absurde. J'en suis sorti, à présent, et une fois dehors ça ne
colle plus au reste, ça ne se raccorde plus. C'est quand je suis seul
– dans la foule, dans le métro – que les souvenirs reprennent leur
consistance. J'étais bien tranquille, bien vide, comme tout le monde,
et tout à coup il y a cette haleine contre mon visage. Je reconnais
l'odeur de cuir et de drap de troupe. J'ai à nouveau la main grasse
sur ma chair. Je redeviens cet homme nu, ses vêtements à ses pieds,
un homme qui a froid, qui a honte de son ventre gonflé et de ses
jambes misérables. Ou bien, c'est le sous-officier allemand qui
surgit. Le vieux sous-officier avec sa veste courte, ses grosses
fesses. Il se tient au bord du trottoir, un bâton à la main, planté
dans ses bottes énormes. Et quand nous passons devant lui, il tape
dans le tas. C'est comme ça qu'ils me tombent dessus, les souvenirs,
qu'ils m'attaquent soudain et pèsent sur moi de leur poids atroce. Ça
ne dure pas. Quelqu'un demande~: Vous descendez à la prochaine~? Les
gens me bousculent, me délivrent.

«~Voilà, c'est un type qui…~» Mon petit récit a du succès. Tout à fait
la sorte de récits qui convient aux familles~: coloré, drôle – et
crâne\footnote{«~Crâne~»~: courageux, décidé.} en même temps~; moitié
Courteline\footnote{Georges Courteline~: dramaturge et romancier
français (1858-1923), dont l'œuvre est essentiellement comique.} et
moitié Déroulède\footnote{Paul Déroulède~: écrivain français
(1846-1914), auteur d'une œuvre empreinte de nationalisme, et
cofondateur de la Ligue des patriotes, organisation
d'extrême-droite.}. La Famille s'amuse et admire. […] Et ainsi, à
mesure que j'en parle, mes cinquante mois de captivité se
transforment en une bonne blague de chambrée, en une partie de
cache-cache avec nos gardiens. Voilà ce que j'aurai rapporté de mon
voyage~: une demi-douzaine d'anecdotes qui feront rigoler la famille
à la fin des repas de famille.

Mes vrais souvenirs, pas question de les sortir.
\theme{guerre!Seconde Guerre mondiale}\theme{camp de concentration}\theme{témoignage}\theme{souvenir}\theme{habitude}\theme{solitude}\theme{violence}\theme{souffrance}\theme{mensonge}\theme{soi!parler de soi}\theme{témoignage!témoigner de la violence}\theme{Shoah}}

%%%%%%%%%%%%%%%%%%%%%%%%%%%%Nv-Sujet%%%%%%%%%%%%%%%%%%%%%%%%%%%%%%%%%%%%%%%%%%

\sujetn{%
	nomauteur=Foucault,
	prenomauteur=Michel,
	initialeprenom =M,
	particule=,
	titre = {Les anormaux},
	semestre = \HQ,
	lieu =Métropole,
	annee =2021,
	type =,
	jour =1,
	datepub =1974-1975,
	ajoutref ={Cours au Collège de France},
	genre =\coursp,
	corrige =,
	liencorrige={met_2021_1_2},
	typequn = \phil,
	typeqdeux = \litt,
	qun = {Qu'est-ce que le monstre permet de comprendre et pourquoi le permet-il~?},
	qdeux ={Dans la littérature et les arts, la figure du monstre ne sert-elle
qu'à faire peur~?},
	themeprogrammequn=HL,
	themeprogrammeqdeux=HL}
{\chapeau{
Le cours au collège de France prononcé par Michel Foucault de janvier
à mars 1975 poursuit les analyses des années précédentes consacrées à
la question du lien entre savoir et pouvoir. Foucault traite ici du
problème des individus dits «~anormaux~» dont la figure du
«~monstre~» est représentative.}

On peut dire que ce qui fait la force et la capacité d'inquiétude du
monstre, c'est que, tout en violant la loi, il la laisse sans voix.
Il piège la loi qu'il est en train d'enfreindre. Au fond, ce que
suscite le monstre, au moment même où par son existence il viole la
loi, ce n'est pas la réponse de la loi elle-même, mais c'est tout
autre chose. Ce sera la violence, ce sera la volonté de suppression
pure et simple, ou encore ce seront les soins médicaux, ou encore ce
sera la pitié. Mais ce n'est pas la loi elle-même, qui répond à cette
attaque que représente pourtant contre elle l'existence du monstre.
Le monstre est une infraction qui se met automatiquement hors la loi,
et c'est là l'une des premières équivoques. La seconde est que le
monstre est, en quelque sorte, la forme spontanée, la forme brutale,
mais, par conséquent, la forme naturelle de la contre-nature. C'est
le modèle grossissant, la forme déployée par les jeux de la nature
elle-même de toutes les petites irrégularités possibles. Et, en ce
sens, on peut dire que le monstre est le grand modèle de tous les
petits écarts. C'est le principe d'intelligibilité de toutes les
formes – circulant sous forme de menue monnaie – de l'anomalie.
Chercher quel est le fond de monstruosité qu'il y a derrière les
petites anomalies, les petites déviances, les petites irrégularités~:
c'est ce problème qui va se retrouver tout au long du XIX\ieme{} siècle.
C'est la question, par exemple, que Lombroso\footnote{Cesare Lombroso
: médecin et spécialiste italien de la criminalité (1835-1909).}
posera lorsqu'il aura affaire à des délinquants. Quel est le grand
monstre naturel qui se profile derrière le petit voleur~? Le monstre
est paradoxalement – malgré la position limite qu'il occupe, bien
qu'il soit à la fois l'impossible et l'interdit – un principe
d'intelligibilité.
\theme{monstre}\theme{politique}\theme{violence}\theme{pitié}\theme{nature!nature humaine}\theme{anomalie}\theme{connaissance}\theme{déportation}}

%%%%%%%%%%%%%%%%%%%%%%%%%%%%Nv-Sujet%%%%%%%%%%%%%%%%%%%%%%%%%%%%%%%%%%%%%%%%%%

\sujetn{%
	nomauteur=Delbo,
	prenomauteur=Charlotte,
	initialeprenom =C,
	particule=,
	titre = {Auschwitz et après},
	semestre = \HQ,
	lieu =Métropole,
	annee =2021,
	type =,
	jour =2,
	datepub =1965,
	ajoutref ={tome I, \textit{Aucun de nous ne reviendra}},
	genre =autobiographie,
	corrige =,
	liencorrige={met_2021_2_1},
	typequn = \litt,
	typeqdeux = \phil,
	qun = {Comment le texte rend-il compte de la violence de la déportation~?},
	qdeux ={À quelles conditions un monde peut-il être dit humain~?},
	themeprogrammequn=HV,
	themeprogrammeqdeux=HVL}
{\vspace*{-24pt}

\chapeau{
Charlotte Delbo, résistante, a été déportée à Auschwitz le 24 janvier
1943. Survivante du camp d'extermination, elle écrit six mois après
son retour de déportation un texte qui ne sera publié que vingt ans
plus tard.}

La gare n'est pas une gare. C'est la fin d'un rail. Ils regardent et
ils sont éprouvés par la désolation autour d'eux.

Le matin la brume leur cache les marais.

Le soir les réflecteurs éclairent les barbelés blancs dans une netteté
de photographie astrale. Ils croient que c'est là qu'on les mène et
ils sont effrayés.

La nuit ils attendent le jour avec les enfants qui pèsent aux bras des
mères. Ils attendent et ils se demandent.

Le jour ils n'attendent pas. Les rangs se mettent en marche tout de
suite. Les femmes avec les enfants d'abord, ce sont les plus
las\footnote{«~las~»~: qui éprouve une fatigue physique, brisé,
éreinté.}. Les hommes ensuite. Ils sont aussi las mais ils sont
soulagés qu'on fasse passer en premier leurs femmes et leurs enfants.
Car on fait passer en premier les femmes et les enfants.

L'hiver ils sont saisis par le froid. Surtout ceux qui viennent de
Candie\footnote{«~Candie~»~: nom médiéval de la ville d'Héraklion, en
Crête.} la neige leur est nouvelle.

L'été le soleil les aveugle au sortir des fourgons obscurs qu'on a
verrouillés au départ.

Au départ de France d'Ukraine d'Albanie de Belgique de Slovaquie
d'Italie de Hongrie du Péloponnèse de Hollande de Macédoine
d'Autriche d'Herzégovine des bords de la mer Noire et des bords de la
Baltique des bords de la Méditerranée et des bords de la Vistule.

Ils voudraient savoir où ils sont. Ils ne savent pas que c'est ici le
centre de l'Europe. Ils cherchent la plaque de la gare. C'est une
gare qui n'a pas de nom. Une gare qui pour eux n'aura jamais de nom.

\bigskip

Il y en a qui voyagent pour la première fois de leur vie. Il y en a
qui ont voyagé dans tous les pays du monde, des commerçants. Tous les
paysages leur étaient familiers mais ils ne reconnaissent pas
celui-ci.

Ils regardent. Ils sauront dire plus tard comme c'était.

Tous veulent se rappeler quelle impression ils ont eue et comme ils
ont eu le sentiment qu'ils ne reviendraient pas.

C'est un sentiment qu'on peut avoir eu déjà dans sa vie. Ils savent qu'il faut se défier des sentiments.
\theme{centre de mise à mort}\theme{camp d'extermination|see{centre de mise à mort}}\theme{déportation}\theme{peur}\theme{violence}\theme{Europe}\theme{déshumanisation}\theme{marche (militaire)}}

%%%%%%%%%%%%%%%%%%%%%%%%%%%%Nv-Sujet%%%%%%%%%%%%%%%%%%%%%%%%%%%%%%%%%%%%%%%%%%
\sujetn{%
	nomauteur=Arendt,
	prenomauteur=Hannah,
	initialeprenom =H,
	particule=,
	titre = {Journal de pensée},
	semestre = \HQ,
	lieu =Métropole,
	annee =2021,
	type =Remplacement,
	jour =1,
	datepub =1953,
	ajoutref ={traduction Sylvie \textsc{Courtine-Denamy}},
	genre =\carnet,
	corrige =,
	liencorrige={metrem_2021_1_2},
	typequn = \phil,
	typeqdeux = \litt,
	qun = {À quelles conditions peut-on, selon Hannah Arendt, prendre le risque
de la guerre~?},
	qdeux ={La littérature et les arts peuvent-ils représenter la guerre sans la dénoncer~?},
	themeprogrammequn=HV,
	themeprogrammeqdeux=HV}
{\chapeau{
Hannah Arendt a tenu toute sa vie un Journal de pensée dans lequel
elle a rédigé des notes diverses qui ont préparé ses livres
philosophiques. En 1953, elle réfléchit notamment à la question de la
guerre.}

\bigskip

À propos de la question de la guerre~: c'est seulement parce que nous
savons que nous devons mourir, par conséquent en mettant les choses
au pis\footnote{«~au pis~»~: au pire}, parce qu'on se dessaisit de quelque chose qui nous sera de
toute façon enlevé, qu'on peut risquer sa vie pour quelque chose. Si
nous étions immortels (non pas comme les dieux qui sont condamnés à
l'être et pour lesquels il n'existe pour cette raison pas de liberté
en général) de telle sorte que nous puissions mourir, mais sans y
être obligés, aucun enjeu au nom duquel on pourrait risquer sa vie ne
serait pensable~: la vie deviendrait simplement un absolu en dehors
duquel il n'y aurait tout bonnement rien. On ne peut sacrifier sa vie
qu'à la liberté parce qu'en dehors de sa propre vie il n'existe pas
de vie du genre humain qui la dépasse. Au cas où l'immortalité de
notre propre vie serait possible, la vie en tant que telle
deviendrait quelque chose d'absolu, au sens où tout ce qu'on appelle
les «~valeurs~» ne pourraient se mouvoir qu'au sein de cette vie.

C'est précisément à cette immortalité possible mais non assurée qu'est
en premier lieu confronté chaque peuple, et en définitive le genre
humain. C'est pourquoi les politiciens nationaux peuvent bien mettre
en jeu la puissance politique et même la liberté politique de leur
peuple, mais jamais toutefois son existence physique elle-même,
puisque de celle-ci dépend précisément le fait qu'il puisse y avoir
en général quelque chose comme une politique. Compte tenu de son
immortalité potentielle, un peuple ne peut jamais être mis en jeu au
nom de quelque chose d'autre. Les limites de toute politique
consistent en ce qu'elles doivent respecter cette possibilité, la
protéger, la garantir. Tout cela est à plus forte raison valable pour
l'humanité. Aucune guerre ne devrait avoir le droit de mettre en jeu
l'existence de l'humanité. Or c'est précisément cela qui est devenu
une possibilité, un risque possible et redouté. La liberté, la
justice, etc., seraient des mots vides s'il s'agissait de la
pérennité\footnote{«~pérennité~»~: permanence} physique de l'humanité ou de la pérennité terrestre de
son habitat, la terre. À l'heure où une destruction de toute vie sur
terre ou la destruction de la terre elle-même n'est pensable que
comme une sorte de «~surprise de la technique~», on ne peut plus
attendre d'aucun peuple qu'il prenne le risque de la guerre.
\theme{guerre}\theme{vie (et mort)}\theme{liberté}\theme{politique}\theme{immortalité}\theme{humanité}\theme{bombe atomique}\theme{vie (et mort)!sens de la vie}\theme{valeurs}}

%%%%%%%%%%%%%%%%%%%%%%%%%%%%Nv-Sujet%%%%%%%%%%%%%%%%%%%%%%%%%%%%%%%%%%%%%%%%%%

\sujetn{%
	nomauteur=Freud,
	prenomauteur=Sigmund,
	initialeprenom =S,
	particule=,
	titre = {Le Malaise dans la civilisation},
	semestre = \HQ,
	lieu =Métropole,
	annee =2021,
	type =Remplacement,
	jour =2,
	datepub =1930,
	ajoutref ={traduit de l'allemand par Bernard \textsc{Lotholary} (2010)},
	genre =\essaip,
	corrige =,
	liencorrige={metrem_2021_2_2},
	typequn = \phil,
	typeqdeux = \litt,
	qun = {D'après ce texte, comment la civilisation peut-elle répondre à
l'agressivité~?},
	qdeux ={La représentation de la violence dans la littérature et les arts
constitue-t-elle un frein à l'agressivité de l'être humain~?},
	themeprogrammequn=HV,
	themeprogrammeqdeux=HV}
{L'existence de cette agressivité que nous pouvons éprouver en
nous-mêmes et supposons à bon droit chez autrui, tel est l'élément
qui perturbe nos rapports avec notre prochain et contraint la
civilisation à tout ce qu'elle met en œuvre. Du fait de cette
hostilité primaire des êtres humains les uns envers les autres, la
société civilisée est constamment menacée de se désagréger. L'intérêt
de la communauté de travail ne la maintiendrait pas soudée, les
passions instinctives sont plus fortes que les intérêts raisonnables.
La civilisation doit tout mettre en œuvre pour dresser des barrières
devant les instincts agressifs des hommes, pour en réduire les
manifestations par des dispositifs psychiques qui réagissent contre.
D'où le recours, donc, à des méthodes visant à inciter les êtres
humains à des identifications et à des relations d'amour
n'aboutissant pas, d'où la restriction imposée à la vie sexuelle et
d'où, également, le commandement idéal d'aimer son prochain comme
soi-même, qui se justifie en réalité par le fait que rien n'est plus
contraire à la nature originelle de l'homme. Quelques peines qu'il
ait coûtées, cet effort de la civilisation n'a pas, jusqu'ici,
remporté grand succès. Les plus grossiers débordements de force
brutale, la civilisation espère y faire obstacle en s'arrogeant le
droit d'exercer elle-même la violence contre les criminels~; mais les
manifestations plus prudentes et plus subtiles de l'agressivité
humaine, la loi n'a pas prise sur elles. Chacun de nous en vient à
laisser tomber comme autant d'illusions les attentes que, dans sa
jeunesse, il rattachait à ses semblables, et chacun peut constater
combien la vie lui est rendue difficile et douloureuse par leur
malignité\footnote{«~malignité~»~: méchanceté.}. Pour autant, il serait injuste de reprocher à la
civilisation de vouloir exclure des actions humaines concurrences et
conflits. Ceux-ci sont sûrement indispensables, mais l'antagonisme
n'est pas nécessairement de l'hostilité, il n'est qu'abusivement pris
pour prétexte de celle-ci.
\theme{violence!psychologie de la violence}\theme{agressivité}\theme{société}\theme{pulsions}\theme{nature!nature humaine}\theme{amour}\theme{civilisation}\theme{culture}\theme{art!religion}}

%%%%%%%%%%%%%%%%%%%%%%%%%%%%Nv-Sujet%%%%%%%%%%%%%%%%%%%%%%%%%%%%%%%%%%%%%%%%%%

\sujetn{%
	nomauteur=Char,
	prenomauteur=René,
	initialeprenom =R,
	particule=,
	titre = {Fureur et mystère},
	semestre = \HQ,
	lieu =Polynésie,
	annee =2021,
	type =,
	jour =1,
	datepub =1948,
	ajoutref ={«~Feuillets d'Hypnos~», fragment 128},
	genre =\autobiographie,
	corrige =,
	liencorrige={pol_2021_1_1},
	typequn = \litt,
	typeqdeux = \phil,
	qun = {D'où vient l'émotion qui se dégage de ce texte~?},
	qdeux ={Comment résister à un agresseur sans recourir à la violence~?},
	themeprogrammequn=HV,
	themeprogrammeqdeux=HV}
{\chapeau{
De 1941 à 1944, René Char tient un rôle actif dans la Résistance. Dans
«~\textup{Feuillets d'Hypnos}~», écrit entre 1943 et 1944, le poète
témoigne de son engagement à travers des fragments poétiques qui
prennent parfois la forme de courts récits.}

\bigskip

Le boulanger n'avait pas encore dégrafé les rideaux de fer de sa
boutique que déjà le village était assiégé, bâillonné, hypnotisé, mis
dans l'impossibilité de bouger. Deux compagnies de SS et un
détachement de miliciens le tenaient sous la gueule de leurs
mitrailleuses et de leurs mortiers. Alors commença l'épreuve.

Les habitants furent jetés hors des maisons et sommés de se rassembler
sur la place centrale. Les clés sur les portes. Un vieux, dur
d'oreille, qui ne tenait pas compte assez vite de l'ordre, vit les
quatre murs et le toit de sa grange voler en morceaux sous l'effet
d'une bombe. Depuis quatre heures j'étais éveillé. Marcelle était
venue à mon volet me chuchoter l'alerte. J'avais reconnu
immédiatement l'inutilité d'essayer de franchir le cordon de
surveillance et de gagner la campagne. Je changeai rapidement de
logis. La maison inhabitée où je me réfugiai autorisait, à toute
extrémité, une résistance armée efficace. Je pouvais suivre de la
fenêtre, derrière les rideaux jaunis, les allées et venues nerveuses
des occupants. Pas un des miens n'était présent au village. Cette
pensée me rassura. À quelques kilomètres de là, ils suivraient mes
consignes et resteraient tapis. Des coups me parvenaient, ponctués
d'injures. Les SS avaient surpris un jeune maçon qui revenait de
relever des collets. Sa frayeur le désigna à leurs tortures. Une voix
se penchait hurlante sur le corps tuméfié~: «~Où est-il~?
Conduis-nous~», suivie de silence. Et coups de pied et coups de
crosse de pleuvoir. Une rage insensée s'empara de moi, chassa mon
angoisse. Mes mains communiquaient à mon arme leur sueur crispée,
exaltaient sa puissance contenue. Je calculais que le malheureux se
tairait encore cinq minutes, puis, fatalement, il parlerait. J'eus
honte de souhaiter sa mort avant cette échéance. Alors apparut
jaillissant de chaque rue la marée des femmes, des enfants, des
vieillards, se rendant au lieu de rassemblement, suivant un
\textit{plan concerté}. Ils se hâtaient sans hâte, ruisselant
littéralement sur les SS, les paralysant «~en toute bonne foi~». Le
maçon fut laissé pour mort. Furieuse, la patrouille se fraya un
chemin à travers la foule et porta ses pas plus loin. Avec une
prudence infinie, maintenant des yeux anxieux et bons regardaient
dans ma direction, passaient comme un jet de lampe sur ma fenêtre. Je
me découvris à moitié et un sourire se détacha de ma pâleur. Je
tenais à ces êtres par mille fils confiants dont pas un ne devait se
rompre.

J'ai aimé farouchement mes semblables cette journée-là, bien au-delà
du sacrifice *.

[* N'était-ce pas le hasard qui m'avait choisi pour prince ce jour-là
plutôt que le cœur mûri pour moi de ce village~? (1945.)]
\theme{guerre!Seconde Guerre mondiale}\theme{violence}\theme{Résistants}\theme{SS}\theme{humanité}\theme{violence!non-violence}\theme{torture}\theme{sacrifice}}

%%%%%%%%%%%%%%%%%%%%%%%%%%%%Nv-Sujet%%%%%%%%%%%%%%%%%%%%%%%%%%%%%%%%%%%%%%%%%%

\sujetn{%
	nomauteur=Echenoz,
	prenomauteur=Jean,
	initialeprenom =J,
	particule=,
	titre = {14},
	semestre = \HQ,
	lieu =Polynésie,
	annee =2021,
	type =Remplacement,
	jour =1,
	datepub =2012,
	ajoutref =chapitre 8,
	genre =\romanl,
	corrige =,
	liencorrige={polrem_2021_1_1},
	typequn = \litt,
	typeqdeux = \phil,
	qun = {Qu'est-ce que la représentation de la guerre a de paradoxal dans cet extrait~?},
	qdeux ={La guerre est-elle toujours absurde~?},
	themeprogrammequn=HV,
	themeprogrammeqdeux=HV}
{\chapeau{\emph{14} évoque le destin de quelques soldats français engagés dans
les combats de la Première Guerre mondiale. Le passage qui suit
constitue leur baptême du feu.}

Puis on leur a crié d'avancer et, plus ou moins poussé par les autres,
il\footnote{Il s'agit d'Anthime, les autres noms propres qui suivent
(Bossis, Arcenel et Padioleau) sont ceux de ses camarades de
régiment.} s'est retrouvé sans trop savoir que faire au milieu d'un
champ de bataille on ne peut plus réel. D'abord avec Bossis ils se
sont regardés, Arcenel derrière eux rajustait une courroie et
Padioleau se mouchait dans un tissu moins blanc que lui. Ensuite il a
bien fallu s'élancer au pas de charge cependant que paraissait à
l'arrière-plan, dans leur dos, un groupe d'une vingtaine d'hommes
qui, le plus paisiblement du monde, se sont disposés en rond sans
apparent souci des projectiles. C'étaient les musiciens du régiment
dont le chef, sa baguette blanche dressée, a fait s'élever en
l'abattant l'air de \textit{La Marseillaise}, l'orchestre envisageant
d'illustrer vaillamment l'assaut. Bien disposés en défense dans un
bois qui les dissimulait, les ennemis ont d'abord empêché la troupe
de progresser mais, l'artillerie s'y mettant par-derrière pour
essayer de les affaiblir, on a entrepris d'attaquer, courant courbés,
maladroitement sous le poids du matériel, chacun précédé de sa
baïonnette qui trouait l'air glacé devant soi.

Or on avait chargé trop tôt, commettant de plus l'erreur de se porter
en masse sur la route qui traversait le théâtre du combat. Cette
route, à découvert et bien repérée par l'artillerie adverse postée
derrière les arbres, constituait en effet une cible parfaitement
dégagée~: tout de suite quelques hommes, pas loin d'Anthime, se sont
mis à tomber, il a cru voir jaillir deux ou trois gerbes de sang mais
les a rejetées avec vigueur de son esprit – n'étant pas même certain,
n'ayant pas le temps d'être certain que ce fût du sang sous pression,
ni d'ailleurs d'en avoir jamais vu à ce jour, du moins pas de cette
façon ni sous cette forme. Il n'avait d'ailleurs pas la tête à
penser, juste à tenter de tirer sur ce qui semblait hostile et,
surtout, chercher un couvert possible où qu'il fût. Par chance,
quoique aussitôt battue en règle par le feu ennemi, la route
présentait çà et là des tronçons encaissés où l'on a d'abord pu
s'abriter un peu.

Mais trop peu~: sous les ordres aboyés, les premiers rangs
d'infanterie ont dû abandonner cette voie pour se risquer ouvertement
dans l'étendue d'avoine qui la bordait et, dès lors, non contents
d'essuyer les tirs venus de l'ennemi, ils ont commencé de recevoir
aussi dans le dos des balles imprudemment tirées par leurs propres
forces, après quoi le désordre s'est vite installé dans les rangs.
C'est qu'on était sans expérience, les accrochages commençaient à
peine~: ce ne serait que plus tard, pour pallier de tels impairs et
se faire mieux repérer par les officiers observateurs, qu'on
recevrait l'ordre de coudre un grand rectangle blanc dans le dos de
sa capote. Cependant, tandis que l'orchestre tenait sa partie dans le
combat, le bras du baryton s'est vu traversé par une balle et le
trombone est tombé, très mauvaisement blessé~: le rond s'est resserré
d'autant et, quoique en formation restreinte, les musiciens ont
continué de jouer sans la moindre fausse note, puis comme ils
reprenaient la mesure où se lève l'étendard sanglant, la flûte et
l'alto sont tombés morts.
\theme{guerre!Première Guerre mondiale}\theme{violence}\theme{guerre}\theme{combat}\theme{bataille}\theme{esthétique!de la violence}\theme{stratégie}\theme{vie (et mort)}\theme{victime}\theme{absurde}\theme{musique}}

%%%%%%%%%%%%%%%%%%%%%%%%%%%%Nv-Sujet%%%%%%%%%%%%%%%%%%%%%%%%%%%%%%%%%%%%%%%%%%

\sujetn{%
	nomauteur=Barrau,
	prenomauteur=Aurélien,
	initialeprenom =A,
	particule=,
	titre = {Le plus grand défi de l'histoire de l'humanité},
	semestre = \HQ,
	lieu ={Amérique du Sud},
	annee =2022,
	type =,
	jour =1,
	datepub =2019,
	ajoutref =,
	genre =\essaip,
	corrige =,
	liencorrige={amsud_2022_1_2},
	typequn = \phil,
	typeqdeux = \litt,
	qun = {Quelle vision de la nature, selon l'auteur, est à l'origine de ce qu'il désigne comme une catastrophe écologique~?},
	qdeux ={Quel rôle vous semble jouer la littérature dans la reconnaissance de l'importance de la nature en elle-même et pour elle-même~?},
	themeprogrammequn=HLHA,
	themeprogrammeqdeux=EXSHL}
{\vspace*{12pt}\\
Par ailleurs, nous savons aujourd'hui que beaucoup d'animaux souffrent comme nous, qu'ils ont une «~conscience~» (au sens le plus fort de ce mot) comme nous, qu'ils ont peur comme nous. Nous le savons et nous les décimons pourtant comme jamais ils ne l'ont été dans l'histoire. Faut-il d'ailleurs absolument qu'ils nous ressemblent pour que nous les aimions et les respections~? Le «~crime contre la vie~» perpétré chaque jour par une humanité plus prédatrice — et de très loin — qu'aucune espèce ne le fut jamais dans l'histoire de la Terre peut-il perdurer indéfiniment~? Allons-nous continuer à l'assumer~? N'est-il pas temps de cesser de faire comme si les vivants non humains étaient des objets alors que nous savons qu'ils ne le sont pas~?

Nous traitons rigoureusement les animaux comme des choses. Unilatéralement, nous avons décidé que la Terre serait l'enfer pour nombre des vivants qui la peuplent. Nous tuons vraisemblablement chaque mois plus d'animaux qu'il n'a existé d'êtres humains dans toute l'Histoire.

La «~nature~» est, elle aussi, souvent pensée sous le seul prisme de ce qu'elle «~rapporte~», de ce qu'elle nous «~dispense~» (comme le furent les peuples colonisés). Peut-être serait-ce l'occasion de la penser pour elle-même. Faut-il continuer à voir les lieux que nous investissons comme étant «~à disposition~»~? Tout un écosystème subtil y préexiste. Il n'est pas une simple ressource. Il ne doit plus être ainsi perçu. Il vaut pour ce qu'il est et non pas pour ce qu'il nous donne. Le problème de la mort massive des animaux et des végétaux est presque toujours présenté du point de vue de ses effets négatifs (souvent réels) sur la vie humaine. Mais n'est-il pas en lui-même catastrophique~? Le monde existe indépendamment de son rôle pour notre confort. La «~loi du plus fort~» n'est pas seulement éthiquement indéfendable, elle se retourne presque toujours contre celui qui en abuse.

Certains pays commencent à donner des droits à des rivières ou à des forêts. D'un point de vue juridique, elles peuvent être représentées de différentes manières (par exemple, par un individu désigné ou par toute personne décidant de porter plainte en cas d'atteinte). C'est une piste intéressante qui mérite d'être explorée. À condition qu'elle ne soit pas une simple poudre aux yeux et que les États aient encore un véritable pouvoir face aux entités supraétatiques qui aujourd'hui gagnent en puissance.

Les parcs nationaux, aussi utiles soient-ils pour sauver quelques bribes de diversité, je l'évoquais précédemment, constituent également l'archétype d'une aberration conceptuelle. Ils soulignent à quel point les humains ont artificiellement décidé que la «~nature~» ne faisait plus partie de leur monde. Elle n'aurait plus droit d'être qu'au sein de sortes de super «~parcs d'attractions~» sous leur contrôle. C'est toute cette logique aberrante qu'il s'agit de renverser.

Le défi que nous avons à relever concerne donc aussi une mutation de nos \emph{valeurs}. L'obligation dans laquelle les pays riches se trouvent de réapprendre un certain «~ascétisme tendanciel~» au niveau matériel n'est pas forcément une mauvaise nouvelle.

\theme{animaux}\theme{conscience}\theme{souffrance}\theme{écologie}\theme{vivant}\theme{nature!environnement}\theme{écosystème}\theme{végétaux}\theme{loi!du plus fort}\theme{droit!des animaux}\theme{État}\theme{justice}\theme{ascétisme}\theme{morale}\theme{humain!et animal}}

%%%%%%%%%%%%%%%%%%%%%%%%%%%%Nv-Sujet%%%%%%%%%%%%%%%%%%%%%%%%%%%%%%%%%%%%%%%%%%
\sujetn{%
	nomauteur=Diop,
	prenomauteur=David,
	initialeprenom =D,
	particule=,
	titre = {Coups de pilon},
	semestre = \HQ,
	lieu ={Amérique du Nord},
	annee =2022,
	type =,
	jour =1,
	datepub =1956,
	ajoutref =,
	genre =\poesie,
	corrige =Oui,
	liencorrige={amnord_2022_1_2},
	typequn = \litt,
	typeqdeux = \phil,
	qun = {Comment ce poème exprime-t-il la violence de l'histoire~?},
	qdeux ={Peut-on surmonter les épreuves de l'histoire~?},
	themeprogrammequn=HV,
	themeprogrammeqdeux=HV}
{\chapeau{
David Diop (1927-1966) est un poète et écrivain français, né en France
de père sénégalais et de mère camerounaise.}

\setlength{\parskip}{0pt}

\ \ \ \ \ \ \ \ Afrique

{\raggedleft
À ma mère
\par}

Afrique, mon Afrique

Afrique des fiers guerriers dans les savanes ancestrales

Afrique que chante ma grand-Mère

Au bord de son fleuve lointain

Je ne t'ai jamais connue

Mais mon regard est plein de ton sang

Ton beau sang noir à travers les champs répandu

Le sang de ta sueur

La sueur de ton travail

Le travail de l'esclavage

L'esclavage de tes enfants

Afrique dis-moi Afrique

Est-ce donc toi ce dos qui se courbe

Et se couche sous le poids de l'humilité

Ce dos tremblant à zébrures rouges

Qui dit oui au fouet sur les routes de midi

Alors gravement une voix me répondit

Fils impétueux cet arbre robuste et jeune

Cet arbre là-bas

Splendidement seul au milieu de fleurs blanches et fanées

C'est l'Afrique ton Afrique qui repousse

Qui repousse patiemment obstinément

Et dont les fruits ont peu à peu

L'amère saveur de la liberté.

\setlength{\parskip}{6pt}
\theme{Afrique}\theme{colonisation}\theme{esclavage}\theme{liberté}\theme{histoire!épreuve de l'histoire}}

%%%%%%%%%%%%%%%%%%%%%%%%%%%%Nv-Sujet%%%%%%%%%%%%%%%%%%%%%%%%%%%%%%%%%%%%%%%%%%

\sujetn{%
	nomauteur=Benjamin,
	prenomauteur=Walter,
	initialeprenom =W,
	particule=,
	titre = {Expérience et pauvreté},
	semestre = \HQ,
	lieu =Amérique du Nord,
	annee =2022,
	type =,
	jour =2,
	datepub =1933,
	ajoutref =trad. Cédric \textsc{Cohen Skalli},
	genre =\essaip,
	corrige =Oui,
	liencorrige={amnord_2022_2_2},
	typequn = \phil,
	typeqdeux = \litt,
	qun = {En quoi, selon Walter Benjamin, l'expérience a-t-elle perdu sa valeur
à la suite de la Première Guerre mondiale~?},
	qdeux ={Toute expérience est-elle communicable par la littérature~?},
	themeprogrammequn=HV,
	themeprogrammeqdeux=EXSHV}
{Dans nos manuels de lecture figurait la fable du vieil homme qui sur
son lit de mort fait croire à ses enfants qu'un trésor est caché dans
sa vigne. Ils n'ont qu'à chercher. Les enfants creusent, mais nulle
trace de trésor. Quand vient l'automne, cependant, la vigne donne
comme aucune autre dans tout le pays. Ils comprennent alors que leur
père a voulu leur léguer le fruit de son expérience~: la vraie
richesse n'est pas dans l'or, mais dans le travail. Ce sont des
expériences de ce type qu'on nous a opposées, en guise de menace ou
d'apaisement, tout au long de notre adolescence~: «~C'est encore
morveux et ça veut donner son avis.~» «~Tu en as encore beaucoup à
apprendre.~» L'expérience, on savait exactement ce que c'était~:
toujours les anciens l'avaient apportée aux plus jeunes. Brièvement,
avec l'autorité de l'âge, sous forme de proverbes~; longuement, avec
sa faconde\footnote{Faconde~: grande facilité de parole.}, sous forme
d'histoires~; parfois dans des récits de pays lointains, au coin du
feu, devant les enfants et les petits-enfants. – Où tout cela est-il
passé~? Trouve-t-on encore des gens capables de raconter une histoire
? Où les mourants prononcent-ils encore des paroles impérissables,
qui se transmettent de génération en génération comme un anneau
ancestral~? Qui, aujourd'hui, sait dénicher le proverbe qui va le
tirer d'embarras~? Qui chercherait à clouer le bec à la jeunesse en
invoquant son expérience passée~?

Non, une chose est claire~: le cours de l'expérience a chuté, et ce
dans une génération qui fit en 1914-1918 l'une des expériences les
plus effroyables de l'histoire universelle. Le fait, pourtant, n'est
peut-être pas aussi étonnant qu'il y paraît. N'a-t-on pas alors
constaté que les gens revenaient muets du champ de bataille~? Non pas
plus riches, mais plus pauvres en expérience communicable. Ce qui
s'est répandu dix ans plus tard dans le flot des livres de guerre
n'avait rien à voir avec une expérience quelconque, car l'expérience
se transmet de bouche à oreille. Non, cette dévalorisation n'avait
rien d'étonnant. Car jamais expériences acquises n'ont été aussi
radicalement démenties que l'expérience stratégique par la guerre de
position, l'expérience économique par l'inflation, l'expérience
corporelle par l'épreuve de la faim, l'expérience morale par les
manœuvres des gouvernants. Une génération qui était encore allée à
l'école en tramway hippomobile\footnote{Tramway hippomobile~: tramway
tiré par des chevaux.} se retrouvait à découvert dans un paysage où
plus rien n'était reconnaissable, hormis les nuages et au milieu,
dans un champ de forces traversé de tensions et d'explosions
destructrices, le minuscule et fragile corps humain.

Cet effroyable déploiement de la technique plongea les hommes dans une
pauvreté tout à fait nouvelle.
\theme{témoignage!témoigner de la violence}\theme{enfant}\theme{éducation}\theme{adolescence}\theme{tradition}\theme{transmission}\theme{guerre!Première Guerre mondiale}\theme{brutalisation}\theme{traumatisme}\theme{technique}}

%%%%%%%%%%%%%%%%%%%%%%%%%%%%Nv-Sujet%%%%%%%%%%%%%%%%%%%%%%%%%%%%%%%%%%%%%%%%%%

\sujetn{%
	nomauteur=Céline,
	prenomauteur=Louis-Ferdinand,
	initialeprenom =L.-F,
	particule=,
	titre = {Voyage au bout de la nuit},
	semestre = \HQ,
	lieu =Asie,
	annee =2022,
	type =,
	jour =1,
	datepub =1932,
	ajoutref =,
	genre =\romanl,
	corrige =,
	liencorrige={asi_2022_1_1},
	typequn = \litt,
	typeqdeux = \phil,
	qun = {Quels sont les effets du travail à l'usine sur le personnage~?},
	qdeux ={Les machines nous font-elles toujours violence~?},
	themeprogrammequn=HV,
	themeprogrammeqdeux=HVL}
{\chapeau{
Dans les années 1930, le narrateur, Bardamu, se fait engager dans les
usines Ford à Detroit, aux États-Unis.}


«~Ça ne vous servira à rien ici vos études, mon garçon~! Vous n'êtes
pas venu ici pour penser, mais pour faire les gestes qu'on vous
commandera d'exécuter… Nous n'avons pas besoin d'imaginatifs dans
notre usine. C'est de chimpanzés dont nous avons besoin… Un conseil
encore. Ne nous parlez plus jamais de votre intelligence~! On pensera
pour vous mon ami~! Tenez-vous-le pour dit.~»

Il avait raison de me prévenir. Valait mieux que je sache à quoi m'en
tenir sur les habitudes de la maison. Des bêtises, j'en avais assez à
mon actif tel quel pour dix ans au moins. Je tenais à passer
désormais pour un petit peinard. Une fois rhabillés, nous fûmes
répartis en files traînardes, par groupes hésitants en renfort vers
ces endroits d'où nous arrivaient les fracas énormes de la mécanique.
Tout tremblait dans l'immense édifice et soi-même des pieds aux
oreilles possédé par le tremblement, il en venait des vitres et du
plancher et de la ferraille, des secousses, vibré de haut en bas. On
en devenait machine aussi soi-même à force et de toute sa viande
encore tremblotante dans ce bruit de rage énorme qui vous prenait le
dedans et le tour de la tête et plus bas vous agitant les tripes et
remontait aux yeux par petits coups précipités, infinis, inlassables.
À mesure qu'on avançait on les perdait les compagnons. On leur
faisait un petit sourire à ceux-là en les quittant comme si tout ce
qui se passait était bien gentil. On ne pouvait plus ni se parler ni
s'entendre. Il en restait à chaque fois trois ou quatre autour d'une
machine.

On résiste tout de même, on a du mal à se dégoûter de sa substance, on
voudrait bien arrêter tout ça pour qu'on y réfléchisse, et entendre
en soi son cœur battre facilement, mais ça ne se peut plus. Ça ne
peut plus finir. Elle est en catastrophe cette infinie boîte aux
aciers et nous on tourne dedans et avec les machines et avec la
terre. Tous ensemble~! Et les mille roulettes et les pilons qui ne
tombent jamais en même temps avec des bruits qui s'écrasent les uns
contre les autres et certains si violents qu'ils déclenchent autour
d'eux comme des espèces de silences qui vous font un peu de bien.
\theme{technique}\theme{machines}\theme{pensée!absence de pensée}\theme{usine}\theme{société!organisation sociale}\theme{bruit}\theme{traumatisme}\theme{silence}\theme{esclavage}}

%%%%%%%%%%%%%%%%%%%%%%%%%%%%Nv-Sujet%%%%%%%%%%%%%%%%%%%%%%%%%%%%%%%%%%%%%%%%%%

\sujetn{%
	nomauteur=Habermas,
	prenomauteur=Jürgen,
	initialeprenom =J,
	particule=,
	titre = {L'intégration républicaine},
	semestre = \HQ,
	lieu =Asie,
	annee =2022,
	type =,
	jour =2,
	datepub =1996,
	ajoutref ={«~Les droits de l'homme. À~l'échelle mondiale et au niveau de l'État~», traduit de l'allemand par\\Rainer \textsc{Rochlitz} (revue)},
	genre =\essaip,
	corrige =,
	liencorrige={asi_2022_2_2},
	typequn = \phil,
	typeqdeux = \litt,
	qun = {Comment Habermas montre-t-il que la paix n'est pas une simple absence de guerre~?},
	qdeux ={Comment la littérature et les autres arts peuvent-ils contribuer à la
paix~?},
	themeprogrammequn=HV,
	themeprogrammeqdeux=HV}
{Pour comprendre les causes complexes de la guerre, il faut concevoir
la paix elle-même comme un processus se déroulant sans intervention
de la force, et qui ne vise pas seulement à empêcher l'emploi de
celle-ci, mais à instaurer les conditions réelles d'une coexistence
sans tensions entre les groupes et entre les peuples. Les
réglementations mises en place ne doivent léser ni l'existence ni la
dignité des intéressés, et ne doivent pas non plus porter atteinte
aux intérêts vitaux et aux sentiments de justice, au point que les
parties en conflit, après avoir épuisé les possibilités qu'offrent
les procédures judiciaires, finissent néanmoins par recourir à la
force. Les programmes politiques qui sont fondés sur un tel concept
de paix auront recours, avant tout emploi de la force militaire, à
tous les moyens disponibles, y compris l'intervention humanitaire,
afin d'agir sur la situation intérieure d'États formellement
souverains pour y favoriser à la fois l'autosuffisance économique et
des conditions sociales supportables, la participation démocratique,
l'établissement de l'État de droit et la tolérance culturelle. De
telles stratégies d'intervention non violente en faveur de processus
de démocratisation tablent sur le fait que, en raison des
interdépendances planétaires, tous les États dépendent aujourd'hui de
leur environnement et sont sensibles au pouvoir des influences
indirectes, exercées «~en douceur~», qui peuvent aller jusqu'à des
sanctions économiques explicites.
\theme{guerre}\theme{paix}\theme{violence!psychologie de la violence}\theme{guerre!causes de la guerre}\theme{État}\theme{État!État de droit}\theme{démocratie}\theme{politique}\theme{droit!droit international}\theme{justice}}


%%%%%%%%%%%%%%%%%%%%%%%%%%%%Nv-Sujet%%%%%%%%%%%%%%%%%%%%%%%%%%%%%%%%%%%%%%%%%%

\sujetn{%
	nomauteur=Weil,
	prenomauteur=Simone,
	initialeprenom =S,
	particule=,
	titre = {Réflexions sur les causes de la liberté et de l'oppression sociale},
	semestre = \HQ,
	lieu =Centres étrangers,
	annee =2022,
	type =,
	jour =1,
	datepub =1934,
	ajoutref =,
	genre =\essaip,
	corrige =Oui,
	liencorrige={cea_2022_1_1},
	typequn = \phil,
	typeqdeux = \litt,
	qun = {Pourquoi Simone Weil en vient-elle à affirmer que «~tout se résume
dans la question du pouvoir~»~?},
	qdeux ={La littérature permet-elle de comprendre la violence dont l'homme est capable~?},
	themeprogrammequn=HV,
	themeprogrammeqdeux=HV}
{Le véritable sujet de \textit{l'Iliade}, c'est l'emprise de la guerre sur les
guerriers, et, par leur intermédiaire, sur tous les humains~; nul ne
sait pourquoi chacun se sacrifie, et sacrifie tous les siens, à une
guerre meurtrière et sans objet, et c'est pourquoi, tout au long du
poème, c'est aux dieux qu'est attribuée l'influence mystérieuse qui
fait échec aux pourparlers de paix, rallume sans cesse les
hostilités, ramène les combattants qu'un éclair de raison pousse à
abandonner la lutte.

Ainsi dans cet antique et merveilleux poème apparaît déjà le mal
essentiel de l'humanité, la substitution des moyens aux fins. Tantôt
la guerre. apparaît au premier plan, tantôt la recherche de la
richesse, tantôt la production~; mais le mal reste le même. Les
moralistes vulgaires se plaignent que l'homme soit mené par son
intérêt personnel~; plût au ciel qu'il en fût ainsi~! L'intérêt est
un principe d'action égoïste, mais borné, raisonnable, qui ne peut
engendrer des maux illimités. La loi de toutes les activités qui
dominent l'existence sociale, c'est au contraire, exception faite
pour les sociétés primitives, que chacun y sacrifie la vie humaine,
en soi et en autrui, à des choses qui ne constituent que des moyens
de mieux vivre. Ce sacrifice revêt des formes diverses, mais tout se
résume dans la question du pouvoir. Le pouvoir, par définition, ne
constitue qu'un moyen~; ou pour mieux dire posséder un pouvoir, cela
consiste simplement à posséder des moyens d'action qui dépassent la
force si restreinte dont un individu dispose par lui-même. Mais la
recherche du pouvoir, du fait même qu'elle est essentiellement
impuissante à se saisir de son objet, exclut toute considération de
fin, et en arrive, par un renversement inévitable, à tenir lieu de
toutes les fins. C'est ce renversement du rapport entre le moyen et
la fin, c'est cette folie fondamentale qui rend compte de tout ce
qu'il y a d'insensé et de sanglant tout au long de l'histoire.
L'histoire humaine n'est que l'histoire de l'asservissement qui fait
des hommes, aussi bien oppresseurs qu'opprimés, le simple jouet des
instruments de domination qu'ils ont fabriqués eux-mêmes, et ravale
ainsi l'humanité vivante à être la chose de choses inertes.
\theme{guerre}\theme{humanité}\theme{religion}\theme{mal}\theme{sacrifice}\theme{esclavage}\theme{intérêt}\theme{histoire}\theme{pouvoir}\theme{valeurs}}

%%%%%%%%%%%%%%%%%%%%%%%%%%%%Nv-Sujet%%%%%%%%%%%%%%%%%%%%%%%%%%%%%%%%%%%%%%%%%%

\sujetn{%
	nomauteur=Freud,
	prenomauteur=Sigmund,
	initialeprenom =S,
	particule=,
	titre = {Considérations actuelles sur la guerre et la mort},
	semestre = \HQ,
	lieu =Centres étrangers,
	annee =2022,
	type =,
	jour =2,
	datepub =1915,
	ajoutref =,
	genre =\essaip,
	corrige =Oui,
	liencorrige={cea_2022_2_2},
	typequn = \phil,
	typeqdeux = \litt,
	qun = {Quels sont les effets de la violence de l'État sur les individus~?},
	qdeux ={Dans quelle mesure la littérature permet-elle de s'interroger sur le rapport de l'individu au pouvoir~?},
	themeprogrammequn=HV,
	themeprogrammeqdeux=HV}
{L'État en guerre se permet toutes les injustices, toutes les
violences, dont la moindre déshonorerait l'individu. Il a recours, à
l'égard de l'ennemi, non seulement à la ruse permise, mais aussi au
mensonge conscient et voulu, et cela dans une mesure qui dépasse tout
ce qui s'était vu dans des guerres antérieures. L'État impose aux
citoyens le maximum d'obéissance et de sacrifices, mais les traite
en mineurs, en leur cachant la vérité et en soumettant toutes les
communications et toutes les expressions d'opinions à une censure qui
rend les gens, déjà déprimés intellectuellement, incapables de
résister à une situation défavorable ou à une sinistre nouvelle. Il
se dégage de tous les traités et de toutes les conventions qui le
liaient à d'autres États, avoue sans crainte sa rapacité et sa soif
de puissance que l'individu doit approuver et sanctionner par
patriotisme.

Qu'on ne vienne pas nous dire que l'État ne peut pas renoncer à avoir
recours à l'injustice, car s'il y renonçait, il se mettrait en état
d'infériorité. Se conformer aux normes morales, renoncer à l'activité
brutale et violente est pour l'individu aussi peu avantageux que pour
l'État, et celui-ci se montre rarement disposé à dédommager le
citoyen des sacrifices qu'il exige de lui. Il ne faut pas, en outre,
s'étonner de constater que le relâchement des rapports moraux entre
les grands individus de l'humanité ait eu ses répercussions sur la
morale privée, car notre conscience, loin d'être le juge implacable
dont parlent les moralistes, est, par ses origines, de l'«~angoisse
sociale~», et rien de plus. Là où le blâme de la part de la
collectivité vient à manquer, la compression des mauvais instincts
cesse, et les hommes se livrent à des actes de cruauté, de perfidie,
de trahison et de brutalité, qu'on aurait crus impossibles, à en
juger uniquement par leur niveau de culture.
\theme{guerre}\theme{État}\theme{justice}\theme{violence}\theme{injustice|see{justice}}\theme{sacrifice}\theme{mensonge}\theme{censure}\theme{patriotisme}\theme{nationalisme|see{patriotisme}}\theme{droit!droit international}\theme{morale!conscience morale}\theme{société!organisation sociale}\theme{pulsions}\theme{culture}\theme{individu}\theme{pouvoir}}

%%%%%%%%%%%%%%%%%%%%%%%%%%%%Nv-Sujet%%%%%%%%%%%%%%%%%%%%%%%%%%%%%%%%%%%%%%%%%%

\sujetn{%
	nomauteur=Lévi-Strauss,
	prenomauteur=Claude,
	initialeprenom =C,
	particule=,
	titre = {Tristes Tropiques},
	semestre = \HQ,
	lieu =Liban,
	annee =2022,
	type =,
	jour =1,
	datepub =1955,
	ajoutref ={«~La fin des voyages~»},
	genre =\essaip,
	corrige ={lib_2022_1_2},
	typequn = \phil,
	typeqdeux = \litt,
	qun = {Selon Lévi-Strauss, voyager, est-ce découvrir de nouveaux mondes ou
est-ce les exploiter~?},
	qdeux ={La littérature et les arts peuvent-ils résister à l'uniformisation des cultures~?},
	themeprogrammequn=HVDMRC,
	themeprogrammeqdeux=HVDMRC}
{\label{lévi-strauss-liban2022}Ce que d'abord vous nous montrez, voyages, c'est notre ordure lancée
au visage de l'humanité.

Je comprends alors la passion, la folie, la duperie des récits de
voyage. Ils apportent l'illusion de ce qui n'existe plus et qui
devrait être encore, pour que nous échappions à l'accablante évidence
que vingt mille ans d'histoire sont joués. Il n'y a plus rien à faire
: la civilisation n'est plus cette fleur fragile qu'on préservait,
qu'on développait, à grand peine, dans quelques coins abrités, un
terroir riche en espèces rustiques, menaçantes sans doute par leur
vivacité, mais qui permettaient aussi de varier et de revigorer les
semis. L'humanité s'installe dans la monoculture~; elle s'apprête à
produire la civilisation en masse, comme la betterave.

On risquait jadis sa vie dans les Indes ou les Amériques pour
rapporter des biens qui nous paraissent aujourd'hui dérisoires~: bois
de braise (d'où Brésil)~: teinture rouge, ou poivre dont, au temps
d'Henri IV, on avait à ce point la folie que la Cour en mettait dans
des bonbonnières des grains à croquer. Ces secousses visuelles ou
olfactives, cette joyeuse chaleur pour les yeux, cette brûlure
exquise pour la langue ajoutaient un nouveau registre au clavier
sensoriel d'une civilisation qui ne s'était pas doutée de sa fadeur.
Dirons-nous alors que par un double renversement, nos modernes Marco
Polo rapportent de ces mêmes terres, cette fois sous forme de
photographies, de livres et de récits, les épices morales dont notre
société éprouve un besoin plus aigu en se sentant sombrer dans
l'ennui~?
\theme{civilisation}\theme{anthropocentrisme}\theme{voyage}\theme{humanité}\theme{nature!nature humaine}\theme{culture!uniformisation des cultures}\theme{divertissement}\theme{colonisation}\theme{commerce}}

%%%%%%%%%%%%%%%%%%%%%%%%%%%%Nv-Sujet%%%%%%%%%%%%%%%%%%%%%%%%%%%%%%%%%%%%%%%%%%

\sujetn{%
	nomauteur=Jonas,
	prenomauteur=Hans,
	initialeprenom =H,
	particule=,
	titre = {Le Principe responsabilité},
	semestre = \HQ,
	lieu =Liban,
	annee =2022,
	type =,
	jour =2,
	datepub =1979,
	ajoutref ={trad. Jean \textsc{Greisch}},
	genre =\essaip,
	corrige =,
	liencorrige={lib_2022_2_2},
	typequn = \phil,
	typeqdeux = \litt,
	qun = {Pourquoi, selon Hans Jonas, la précaution vaut-elle mieux que
l'espérance~?},
	qdeux ={Dans un monde dominé par la technique, la littérature nous aide-elle à nous orienter~?},
	themeprogrammequn=HL,
	themeprogrammeqdeux=HL}
{L'expérience a prouvé que les développements déclenchés à chaque fois
par l'agir technologique afin de réaliser des buts à court terme ont
tendance à se rendre autonomes, c'est-à-dire à acquérir leur propre
dynamique contraignante, une inertie autonome, en vertu de laquelle
ils ne sont pas seulement irréversibles, comme on l'a déjà dit, mais
qu'ils poussent également en avant et qu'ils débordent le vouloir et
la planification de ceux qui agissent. Ce qui a été commencé nous ôte
l'initiative de l'agir et les faits accomplis que le commencement a
créés s'accumulent pour devenir la loi de sa continuation. Même s'il
se peut que «~nous prenions en main notre propre développement~»,
celui-ci échappera à nos mains simplement du fait qu'il s'est
incorporé son impulsion et plus que partout ailleurs vaut ici la loi
qu'alors que le premier pas relève de notre liberté, nous sommes
esclaves du second et de tous ceux qui suivent. Ainsi au constat que
l'accélération du développement alimenté technologiquement ne laisse
plus le temps pour des corrections automatiques s'ajoute le constat
ultérieur que pendant le temps que malgré tout nous avons encore à
notre disposition, la correction devient de plus en plus difficile et
la liberté pour la faire diminue continuellement. Cela renforce
l'obligation de veiller aux commencements, accordant la priorité aux
possibilités de malheur fondées de manière suffisamment sérieuse (et
distinctes des simples fantasmes de la peur) par rapport aux
espérances – même si celles-ci ne sont pas moins bien fondées.
\theme{technique!maîtrise technique de la nature}\theme{technologie}\theme{liberté}\theme{écologie}\theme{esclavage}\theme{peur}\theme{espérances}\theme{principe de précaution}}

%%%%%%%%%%%%%%%%%%%%%%%%%%%%Nv-Sujet%%%%%%%%%%%%%%%%%%%%%%%%%%%%%%%%%%%%%%%%%%

\sujetn{%
	nomauteur=Ricœur,
	prenomauteur=Paul,
	initialeprenom =P,
	particule=,
	titre = {Histoire et vérité},
	semestre = \HQ,
	lieu =Métropole,
	annee =2022,
	type =,
	jour =1,
	datepub =1955,
	ajoutref =,
	genre =\essaip,
	corrige =Oui,
	liencorrige={met_2022_1_2},
	typequn = \phil,
	typeqdeux = \litt,
	qun = {D'après l'auteur, qu'est-ce qui explique la permanence de la violence
dans l'histoire~?},
	qdeux ={La littérature et les arts naissent-ils de «~l'appétit de
catastrophe~» des hommes~?},
	themeprogrammequn=HV,
	themeprogrammeqdeux=HV}
{Que la violence soit de toujours et de partout, il n'est que de
regarder comment s'édifient et s'écroulent les empires, s'installent
les prestiges personnels, s'entre-déchirent les religions, se
perpétuent et se déplacent les privilèges de la propriété et du
pouvoir, comment même se consolide l'autorité des maîtres à penser,
comment se juchent\footnote{«~se juchent~»~: se hissent} les
jouissances culturelles des élites sur le tas des travaux et des
douleurs des déshérités.

On ne voit jamais assez grand quand on prospecte l'empire de la
violence~; c'est pourquoi une anatomie de la guerre qui se flatterait
d'avoir découvert trois ou quatre grosses ficelles qu'il suffirait de
couper pour que les marionnettes militaires retombent inertes sur les
tréteaux condamnerait le pacifisme à rester superficiel et puéril.
Une anatomie de la guerre requiert la tâche plus vaste d'une
physiologie de la violence.

Il faudrait aller chercher très bas et très haut les complicités d'une
affectivité humaine accordée au terrible dans l'histoire. La
psychologie sommaire de l'empirisme qui gravite autour du plaisir et
de la douleur, du bien-être et du bonheur, omet
l'irascible\footnote{«~irascible~»~: qui se met facilement en
colère}, le goût de l'obstacle, la volonté d'expansion, de combat et
de domination, les instincts de mort et surtout cette capacité de
destruction, cet appétit de catastrophe qui est la contrepartie de
toutes les disciplines qui font de l'édifice psychique de l'homme un
équilibre instable et toujours menacé. Que l'émeute explose dans la
rue, que la patrie soit proclamée en danger, quelque chose en moi est
rejoint et délié, à quoi ni le métier, ni le foyer, ni les
quotidiennes tâches civiques ne donnaient issue~; quelque chose de
sauvage, quelque chose de sain et de malsain, de jeune et d'informe,
un sens de l'insolite, de l'aventure, de la disponibilité, un goût
pour la rude fraternité et pour l'action expéditive, sans médiation
juridique et administrative. L'admirable est que ces dessous de la
conscience resurgissent au niveau des plus hautes couches de la
conscience~: ce sens du terrible est aussi le sens idéologique~;
soudain la justice, le droit, la vérité prennent des majuscules en
prenant les armes et en s'auréolant de sombres passions.
\theme{violence!psychologie de la violence}\theme{guerre!causes de la guerre}\theme{religion}\theme{pouvoir}\theme{pulsions}\theme{dialectique}\theme{valeurs}\theme{inconscient}\theme{droit}\theme{idéologie}\theme{esthétique!de la violence}}

%%%%%%%%%%%%%%%%%%%%%%%%%%%%Nv-Sujet%%%%%%%%%%%%%%%%%%%%%%%%%%%%%%%%%%%%%%%%%%

\sujetn{%
	nomauteur=Calaferte,
	prenomauteur=Louis,
	initialeprenom =L,
	particule=,
	titre = {C'est la guerre},
	semestre = \HQ,
	lieu =Métropole,
	annee =2022,
	type =,
	jour =2,
	datepub =1993,
	ajoutref =,
	genre =\romanl,
	corrige =Oui,
	liencorrige={met_2022_2_2},
	typequn = \litt,
	typeqdeux = \phil,
	qun = {«~C'est la guerre~»~: ce texte vous en donne-t-il l'impression~?},
	qdeux ={Qu'est-ce qu'être en guerre~?},
	themeprogrammequn=HV,
	themeprogrammeqdeux=HV}
{\chapeau{Ce récit s'ouvre sur l'annonce de la mobilisation générale en
septembre 1939.}

Il est cinq heures d'un après-midi de septembre tiède et gris.

Le tocsin sonne.

On arrête de jouer.


\bigskip

Robe noire fermée jusqu'au cou, les bras levés, des mains blanches
osseuses, le regard fixe, la vieille femme crie sur la place du
village que c'est la mobilisation générale.

Il n'y a pas un souffle d'air dans les feuilles du gros arbre.

Des oiseaux chantent.

\bigskip

Au garde-à-vous dans sa salopette de travail, les mains dans les
poches, un homme pleure.

Il est en sabots.

Il y a du bruit et du silence, mais le silence absorbe le bruit. C'est
comme aux enterrements.

Un long chat noir est étiré sur le rebord d'une fenêtre.

Un long chat noir est étiré sur le rebord d'une fenêtre.

Deux femmes âgées s'étreignent, chacune la tête dans le cou de
l'autre. Le chignon de la plus petite s'est défait, ses cheveux
grisonnants tombent en longues mèches ondulantes de chaque côté de
ses épaules. On dirait des anguilles vivantes. J'ai envie de faire
pipi.

Quelque part, au loin, une génisse appelle d'un meuglement plaintif.

Des villageois restent adossés à la façade jaune sale d'une maison.

Assise sur une pierre, la petite fille bleue tient à deux mains son
ballon sur ses genoux. Ses chaussettes blanches sont en boules molles
sur ses chevilles. Elle se mord les lèvres.

Devant le muret de pierres sèches, une femme s'est agenouillée sur le
sable de la place. Elle a les mains jointes, le dos voûté, la tête
baissée. C'est comme une statue d'église, mais noire.

Ma culotte est trop courte, elle me tire entre les jambes, j'ai de
grosses croûtes aux genoux, ça sanguinole toujours un peu et ça
brûle.

En blouses grises, l'épicier et sa femme se tiennent sur le pas de
leur porte.

Un cerf-volant rouge clignote dans le ciel.

Des hommes arrivent. Ils se serrent la main. On les voit se parler,
hocher la tête, la secouer, hausser les épaules.

Les bras ballants, deux femmes ont déposé devant elles leurs seaux de
fer pleins d'eau.

Je n'ai pas goûté. J'ai faim.

Le petit rouquin se traîne à quatre pattes dans la poussière en
faisant des bulles de salive avec ses lèvres. Il reçoit un coup de
pied, tombe en avant sur le ventre et éclate de rire. C'est sa mère
qui lui a donné le coup de pied. Elle le relève en le tirant
brutalement par le bras. Elle époussette du bout des doigts son
tablier d'écolier noir. Elle lui donne une gifle. Il pleure.


\bigskip

—~On ne tape pas les petits aujourd'hui, dit un vieux, c'est la
guerre.


\bigskip

Je ne sais pas ce que c'est que la mobilisation générale, mais je suis
bien content que ce soit la guerre.

J'ai onze ans.


\bigskip

—~Les salauds dit un homme.

J'aime les tartines épaisses avec dessus du beurre salé et un sucre.

Une grande femme surgit soudainement.

—~Je le savais~! Je le savais~!

Ses cheveux courts semblent grésiller autour de sa tête.

—~Ce matin j'ai écrit une lettre à quelqu'un. Au lieu de mettre la
bonne date j'ai mis deux fois 1914~!

Je la regarde, étroite, nerveuse, les yeux écarquillés, cette voix
criarde. Je ne comprends pas ce qu'elle est en train de dire, mais je
la trouve bête.

—~Papa a fait 14~!

—~Mon père aussi, dit un jeune paysan, le torse nu avec des poils
blonds.

—~Et nous voilà bons encore une fois dit l'homme à la moustache.

Il faut que j'aille chercher mon goûter à la maison.
\theme{enfant}\theme{guerre!Seconde Guerre mondiale}\theme{mobilisation}\theme{absurde}\theme{incompréhension}\theme{tristesse}\theme{patriotisme}\theme{esthétique!de la violence}}

%%%%%%%%%%%%%%%%%%%%%%%%%%%%%Nv-Sujet%%%%%%%%%%%%%%%%%%%%%%%%%%%%%%%%%%%%%%%%

\sujetn{%
	nomauteur=Arendt,
	prenomauteur=Hannah,
	initialeprenom =H,
	particule=,
	titre = {Du mensonge à la violence},
	semestre = \HQ,
	lieu =Métropole,
	annee =2022,
	type =Remplacement,
	jour =1,
	datepub =1972,
	ajoutref =– traduit de l'anglais par Guy \textsc{Durand},
	genre =\essaip,
	corrige =Oui,
	liencorrige={metrem_2022_1_1},
	typequn = \phil,
	typeqdeux = \litt,
	qun = {Pourquoi Hannah Arendt distingue-t-elle le criminel de celui qui fait acte de désobéissance civile~?},
	qdeux ={Écrire est-il un acte de désobéissance~?},
	themeprogrammequn=HV,
	themeprogrammeqdeux=HV}
{La désobéissance civile ne saurait être assimilée à la délinquance de droit commun\footnote{Le délinquant de droit commun est celui qui commet un délit ou un crime, au regard des règles établies du droit.}. Il existe une différence essentielle entre le criminel qui prend soin de dissimuler à tous les regards ses actes répréhensibles et celui qui fait acte de désobéissance civile en défiant les autorités et s'institue lui-même porteur d'un autre droit. Cette distinction nécessaire entre une violation ouverte et publique de la loi et une violation clandestine a un tel caractère d'évidence que le refus d'en tenir compte ne saurait provenir que d'un préjugé allié à de la mauvaise volonté. Reconnue désormais par tous les auteurs sérieux qui abordent ce sujet, cette distinction est naturellement invoquée comme un argument primordial par tous ceux qui s'efforcent de faire reconnaître que la désobéissance civile n'est pas incompatible avec les lois et les institutions publiques des États-Unis. Le délinquant de droit commun, par contre, même s'il appartient à une organisation criminelle, agit uniquement dans son propre intérêt~; il refuse de s'incliner devant la volonté du groupe, et ne cédera qu'à la violence des services chargés d'imposer le respect de la loi. Celui qui fait acte de désobéissance civile, tout en étant généralement en désaccord avec une majorité, agit au nom et en faveur d'un groupe particulier. Il lance un défi aux lois et à l'autorité établie à partir d'un désaccord fondamental, et non parce qu'il entend personnellement bénéficier d'un passe-droit.
\theme{désobéissance}\theme{délinquance}\theme{droit}\theme{autorité}\theme{public}\theme{crime}\theme{police}\theme{écrire}\theme{écriture}}

%%%%%%%%%%%%%%%%%%%%%%%%%%%%%Nv-Sujet%%%%%%%%%%%%%%%%%%%%%%%%%%%%%%%%%%%%%%%%

\sujetn{%
	nomauteur=Levinas,
	prenomauteur=Emmanuel,
	initialeprenom =E,
	particule=,
	titre = {Totalité et infini},
	semestre = \HQ,
	lieu =Métropole,
	annee =2022,
	type =Remplacement,
	jour =2,
	datepub =1974,
	ajoutref =,
	genre =\essaip,
	corrige =Oui,
	liencorrige={metrem_2022_2_1},
	typequn = \phil,
	typeqdeux = \litt,
	qun = {D'après Levinas, en quoi la haine est-elle insatiable~?},
	qdeux ={Pourquoi les sentiments inhumains occupent-ils tant de place dans la littérature et les arts~?},
	themeprogrammequn=HV,
	themeprogrammeqdeux=EXSHV}
{L'épreuve suprême de la liberté – n'est pas la mort mais la souffrance. La haine le sait fort bien qui cherche à saisir l'insaisissable, à humilier, de très haut, à travers la souffrance où autrui existe comme pure passivité~; mais la haine veut cette passivité dans l'être éminemment actif qui doit en témoigner. La haine ne désire pas toujours la mort d'autrui ou, du moins, elle ne désire la mort d'autrui qu'en infligeant cette mort comme une suprême souffrance. Le haineux cherche à être cause d'une souffrance dont l'être haï doit être témoin. Faire souffrir, ce n'est pas réduire autrui au rang d'objet, mais au contraire le maintenir superbement dans sa subjectivité. Il faut que dans la souffrance le sujet sache sa réification\footnote{«~réification~»~: le fait de traiter quelqu'un comme une chose.}, mais pour cela il faut précisément que le sujet demeure sujet. Le haineux veut les deux. D'où le caractère insatiable de la haine~; elle est satisfaite précisément lorsqu'elle ne l'est pas, puisqu'autrui ne la satisfait qu'en devenant objet, mais il ne saurait devenir jamais assez objet puisqu'on exige, en même temps que sa déchéance, sa lucidité et son témoignage. Là réside l'absurdité logique de la haine.
\theme{souffrance}\theme{vulnérabilité}\theme{haine}\theme{vie (et mort)}\theme{autrui}\theme{sujet}\theme{liberté}}

%%%%%%%%%%%%%%%%%%%%%%%%%%%%Nv-Sujet%%%%%%%%%%%%%%%%%%%%%%%%%%%%%%%%%%%%%%%%%%

\sujetn{%
	nomauteur=Genevoix,
	prenomauteur=Maurice,
	particule=,
	initialeprenom =M,
	titre = {Ceux de 14},
	semestre = \HQ,
	lieu =Nouvelle-Calédonie,
	annee =2022,
	type =,
	jour =1,
	datepub =1949,
	ajoutref =,
	genre =\autobiographie,
	corrige =,
	liencorrige={nca_2022_1_1},
	typequn = \litt,
	typeqdeux = \phil,
	qun = {En quoi cette rencontre est-elle particulièrement expressive de la violence humaine~?},
	qdeux ={La violence nous rend-elle insensibles~?},
	themeprogrammequn=HV,
	themeprogrammeqdeux=EXSHV}
{\chapeau{Maurice Genevoix, né en 1880, raconte ses propres souvenirs de la guerre de 14-18. Mobilisé dès le début du conflit, il est affecté à un des secteurs les plus difficiles, celui des Éparges, dans le Nord-Est de la France. Alors qu'il vient d'échapper à un obus, à proximité du front, il voit arriver un cheval.}

Il s'est arrêté court dès qu'il m'a vu. Il reste là, immobile sur ses pattes enflées, les naseaux battants, une oreille pointée vers moi, l'autre tendue en arrière, du côté où sifflaient les balles. Mais bientôt son col s'incline au poids de sa grosse tête et, le mufle à ras de terre, la lèvre longue, il se met à tondre l'herbe.

«~On est ami, n'est-ce pas~?~»

Je caresse le flanc décharné, la peau tiède tendue sur les cercles de la carcasse.

«~Tu saignes, mon pauvre vieux~? Est-ce qu'ils t'auraient touché~?~»

Un filet vermeil sinue au poitrail, glisse le long de la patte gauche, jusqu'au genou. Cela coule d'un petit sillon sombre, près de l'épaule, creusé au passage par la pointe d'une balle.

«~Eh bien~! tu l'as échappé belle~! C'est idiot de se promener comme ça au nez des Boches~!~»

Le vieux cheval a soulevé la tête, dressé l'oreille comme s'il m'écoutait. Mais ses naseaux s'ouvrent tout grands et ses jambes se mettent à trembler~: un obus siffle au loin, franchit la vallée en ronronnant, et plante une colonne de fumée jaune au-dessous du Bois-Haut, à mi-pente. Lorsque le roulement de l'explosion passe sur nous, la pauvre bête, d'un saut maladroit, fait volte-face pour fuir. Plus agile qu'elle, je lui ai barré la route de mes deux bras étendus~: elle recule peu à peu devant moi, la tête rejetée en arrière, ses sabots faisant rouler les pierres. Quand je la vois calmée, retombée à sa placidité, je cours glaner dans une grange quelques poignées de foin perdues aux coins de l'aire. Je reviens. Il est toujours là, broutant à petits coups de lèvres.

«~Tiens, mon bonhomme, c'est pour toi. Mais il faut venir chercher ça de l'autre côté des maisons. Si tu restes par ici, tu vas retourner dans les champs~; et ils te tueront.~» Les grands yeux troubles me regardent, voilés parfois d'un lent clignement. Il flotte dans leur eau profonde un infini de stupeur triste. «~Oui, je comprends~: tu es un vieux cheval très las. L'abri que te donnaient tes maîtres, chaque soir, en récompense de ton labeur du jour, tu ne l'as plus, ni le râtelier plein de foin, ni la musette gonflée d'avoine. Tu es devenu si maigre que tes os crèvent ta peau. Tu as eu si peur, tant de fois, que tes genoux ne cessent de trembler. Et cela dure. Et qu'à la fin ceux de là-bas te tuent, cela, n'est-ce pas, t'est bien égal~?…~»
\theme{insensibilité}\theme{animaux}\theme{guerre}\theme{bataille}\theme{armes}\theme{combat}\theme{amitié}\theme{souffrance}\theme{blessure}}

%%%%%%%%%%%%%%%%%%%%%%%%%%%%Nv-Sujet%%%%%%%%%%%%%%%%%%%%%%%%%%%%%%%%%%%%%%%%%%

\sujetn{%
	nomauteur=Gusdorf,
	prenomauteur=Georges,
	initialeprenom =G,
	particule=,
	titre = {La Vertu de force},
	semestre = \HQ,
	lieu ={Nouvelle-Calédonie},
	annee =2022,
	type =,
	jour =1,
	datepub =2000,
	ajoutref =,
	genre =\essaip,
	corrige =,
	liencorrige={nca_2022_1_2},
	typequn = \phil,
	typeqdeux = \litt,
	qun = {Selon Gusdorf, sommes-nous maîtres ou victimes de notre propre violence~?},
	qdeux={Que nous révèlent de la violence les représentations littéraires et artistiques que nous en avons~?},
	themeprogrammequn=MMHV,
	themeprogrammeqdeux=EXSHV}
{Toute violence, par-delà le meurtre du prochain, poursuit son propre suicide. Elle est en effet destruction de soi~; les Anciens savaient déjà que la colère est une courte folie. La violence suppose un échappement au contrôle~: l'explosion émotive se libère en déchaînements paroxystiques, cris et gesticulations, qui attestent l'échec de toutes les disciplines personnelles. Le violent, incapable de se contenir, recherche dans sa propre frénésie, une sorte d'apaisement magique, comme si, en augmentant le volume et l'intensité de sa voix, en enflant les muscles, il retrouvait cette majorité qu'il sent, devant l'obstacle, confusément perdue. La décharge affective et musculaire peut au surplus procurer le retour au calme, et d'ailleurs le regret de l'excès commis, la honte pour s'être conduit comme un enfant.

Mais il arrive que le violent, une fois hors de soi, ne puisse à nouveau se posséder. Il fait confiance à la violence, méthodiquement, comme on le voit dans le domaine de la terreur, instrument jadis et naguère, et aujourd'hui encore, de la fausse certitude. La violence se fait institution et moyen de gouvernement~: dragonnades\footnote{Les dragonnades sont des persécutions contre les communautés protestantes du royaume de France durant les années 1680, sous le règne de Louis XIV.}, inquisition, univers concentrationnaire et régimes policiers~; il a existé, il existe une civilisation de la violence, monstrueuse affirmation de la certitude qui rend fou, selon la parole de Nietzsche. À travers l'histoire, les persécutions et les guerres maintiennent le pire témoignage que l'humanité puisse porter contre elle-même. Individuelle ou collective, cette violence n'est d'ailleurs que le camouflage d'une faiblesse ressentie, d'un effroi de soi à soi, que l'on essaie, par tous les moyens, de dissimuler.
\theme{violence}\theme{meurtre}\theme{soi}\theme{destruction}\theme{émotions}\theme{soi!contrôle de soi}\theme{corps}\theme{terreur}\theme{camp de concentration}\theme{police}\theme{civilisation}\theme{humanité}\theme{nature!nature humaine}\theme{faiblesse}}

%%%%%%%%%%%%%%%%%%%%%%%%%%%%Nv-Sujet%%%%%%%%%%%%%%%%%%%%%%%%%%%%%%%%%%%%%%%%%%

\sujetn{%
	nomauteur=Jaspers,
	prenomauteur=Karl,
	initialeprenom =K,
	particule=,
	titre = {La bombe atomique et l'avenir de l'homme},
	semestre = \HQ,
	lieu =Nouvelle-Calédonie,
	annee =2022,
	type =,
	jour =2,
	datepub =1963,
	ajoutref ={trad. Edmond \textsc{Saget}},
	genre =\essaip,
	corrige =,
	liencorrige={nca_2022_2_2},
	typequn = \phil,
	typeqdeux = \litt,
	qun = {Quelle est selon Jaspers la contradiction à laquelle aboutit la technique~?},
	qdeux ={Comment la littérature et les arts interrogent-ils le progrès technique~?},
	themeprogrammequn=HL,
	themeprogrammeqdeux=HLCC}
{L'éternel cheminement de la technique, sans lequel le chemin parcouru par l'homme est impensable, est parvenu aujourd'hui en un point où la question se pose à nouveau pour l'homme de savoir ce qu'il veut faire de cette technique. De tout temps, la technique a servi à donner une forme constructive au monde qui nous entoure, de tout temps aussi elle a servi à détruire. Aujourd'hui les possibilités qu'elle offre ont fait le saut qui mène des destructions de détail à la destruction totale de toute vie à la surface de la terre.

La peur de ce danger a conduit à penser qu'il vaudrait mieux n'avoir jamais trouvé la libération de l'énergie atomique, par conséquent n'avoir jamais non plus construit la bombe atomique. Si nous pouvions choisir, il nous faudrait renoncer à l'énergie atomique, accepter plutôt toutes les difficultés qui résultent de la limitation de l'énergie dont nous disposons par ailleurs, que nous exposer à ce péril. Penser ainsi, c'est conclure nécessairement~: il eût mieux valu renoncer au développement technique en général. Car ce développement, une fois mis en marche, ne se laisse pas stopper à un endroit donné ou même ramener en arrière, si ce n'est en détruisant la vie, destruction qui, en mettant fin à la vie du porteur de la technique, mettrait aussi un terme à celle-ci. Le refus d'admettre le dernier pas de la technique a pour conséquence le refus du commencement de la technique.
\theme{progrès}\theme{technique}\theme{destruction}\theme{bombe atomique}\theme{énergie}\theme{électricité}\theme{technique!et sciences}\theme{vie (et mort)}\theme{technique!maîtrise technique de la nature}}

%%%%%%%%%%%%%%%%%%%%%%%%%%%%Nv-Sujet%%%%%%%%%%%%%%%%%%%%%%%%%%%%%%%%%%%%%%%%%%

\sujetn{%
	nomauteur=Arendt,
	prenomauteur=Hannah,
	initialeprenom =H,
	particule=,
	titre = {La condition de l'homme moderne},
	semestre = \HQ,
	lieu =Polynésie,
	annee =2022,
	type =,
	jour =1,
	datepub =1958,
	ajoutref ={trad. Georges \textsc{Fradier}},
	genre =\essaip,
	corrige =,
	liencorrige={pol_2022_1_2},
	typequn = \phil,
	typeqdeux = \litt,
	qun = {Selon Hannah Arendt, est-il souhaitable pour l'être humain d'échapper à sa condition~?},
	qdeux ={Quels sont les pouvoirs de la littérature pour donner sens à la vie humaine~?},
	themeprogrammequn=HL,
	themeprogrammeqdeux=OdH}
{L'artifice humain du monde sépare l'existence humaine de tout milieu
purement animal, mais la vie elle-même est en dehors de ce monde
artificiel, et par la vie l'homme demeure lié à tous les autres
organismes vivants. Depuis quelque temps~; un grand nombre de
recherches scientifiques s'efforcent de rendre la vie
«~artificielle~» elle aussi, et de couper le dernier lien qui
maintient encore l'homme parmi les enfants de la nature. C'est le
même désir d'échapper à l'emprisonnement terrestre qui se manifeste
dans les essais de création en éprouvette, dans le vœu de combiner
«~au microscope le plasma germinal provenant de personnes aux
qualités garanties, afin de produire des êtres supérieurs~» et «~de
modifier (leurs) tailles, formes et fonctions~»~; et je soupçonne que
l'envie d'échapper à la condition humaine expliquerait aussi l'espoir
de prolonger la durée de l'existence fort au-delà de cent ans, limite
jusqu'ici admise.

Cet homme futur, que les savants produiront, nous disent-ils, en un
siècle, pas davantage, paraît en proie à la révolte contre
l'existence humaine telle qu'elle est donnée, cadeau venu de nulle
part (laïquement parlant) et qu'il veut pour ainsi dire échanger
contre un ouvrage de ses propres mains. Il n'y a pas de raison de
douter que nous soyons capables de faire cet échange, de même qu'il
n'y a pas de raison de douter que nous soyons capables à présent de
détruire toute vie organique sur terre. La seule question est de
savoir si nous souhaitons employer dans ce sens nos nouvelles
connaissances scientifiques et techniques, et l'on ne saurait en
décider par des méthodes scientifiques. C'est une question politique
primordiale que l'on ne peut guère, par conséquent, abandonner aux
professions de la science ni à ceux de la politique.
\theme{vie (et mort)}\theme{naturel (et artificiel)}\theme{eugénisme}\theme{biogénétique}\theme{transhumanisme}\theme{technique!maîtrise technique de la nature}\theme{technique!et sciences}\theme{immortalité}\theme{vie (et mort)!sens de la vie}\theme{nature!nature humaine}}

%%%%%%%%%%%%%%%%%%%%%%%%%%%%Nv-Sujet%%%%%%%%%%%%%%%%%%%%%%%%%%%%%%%%%%%%%%%%%%

\sujetn{%
	nomauteur=Martin du Gard,
	prenomauteur=Roger,
	initialeprenom =R,
	particule=,
	titre = {Les Thibault},
	semestre = \HQ,
	lieu =Polynésie,
	annee =2022,
	type =,
	jour =2,
	datepub =1953,
	ajoutref =,
	genre =\romanl,
	corrige =,
	liencorrige={pol_2022_2_1},
	typequn = \litt,
	typeqdeux = \phil,
	qun = {Quelle vision de l'humanité Antoine développe-t-il dans cet extrait~?},
	qdeux ={Peut-on envisager un monde sans guerre~?},
	themeprogrammequn=HV,
	themeprogrammeqdeux=HV}
{\chapeau{À la fin de la Première Guerre mondiale, Antoine est soigné dans une
clinique militaire. Lors d'une permission, il retrouve Daniel mutilé
à la suite d'une blessure.}

Antoine, à son tour, commençait, d'une voix hésitante, entrecoupée de
toux~:

«~Il y a des moments où je me dis que c'est la dernière~; que, après
celle-là, non , il n'est pas possible de pense qu'il puisse y en
avoir d'autres~!… Des moments, où j'en suis sûr… Mais, à d'autres
moments, je doute… Je ne sais plus…~»

Daniel mastiquait en silence, les regards perdus. Que pensait-il~?

Antoine s'était tu. Il avait trop de peine à parler plusieurs minutes
de suite. Mais il continuait à réfléchir aux mêmes choses, pour la
centième, pour la millième fois. «~On est épouvanté, se disait-il,
quand on mesure froidement tout ce qui s'oppose à la pacification
entre les hommes… Combien de siècles encore avant que l'évolution
morale – s'il y a une évolution morale~? – ait enfin purgé l'humanité
de son intolérance instinctive, de son respect inné de la force
brutale, de ce plaisir fanatique qu'éprouve l'animal humain à
triompher par la violence, à imposer, par la violence, ses façons de
sentir, de vivre, à ceux, plus faibles, qui ne sentent pas , qui ne
vivent pas, comme lui~?… Et puis, il y a la politique, les
gouvernements… Pour l'autorité qui déclenche la guerre, pour les
hommes au pouvoir qui la décident et la font faire aux autres, ce
sera toujours, aux heures de faillite, une solution si tentante, si
facile… Peut-on espérer que jamais plus les gouvernements n'y auront
recours~?… Il faudrait alors que ce leur soit devenu impossible~: il
faudrait que le pacifisme ait de telles racines dans l'opinion, ait
pris une telle extension, qu'il oppose un infranchissable obstacle à
la politique belliqueuse des États. C'est chimère que d'espérer ça…
Et puis, le triomphe du pacifisme serait-il seulement une sérieuse
garantie de paix~? Même si, un jour, dans nos pays, les partis
pacifistes tenaient le pouvoir, qui nous dit qu'ils ne céderaient pas
à la tentation de faire de faire la guerre pour le plaisir d'imposer,
par la violence, l'idéologie pacifiste au reste du monde~?…~»
\theme{guerre}\theme{blessure}\theme{paix}\theme{morale}\theme{humanité}\theme{violence!psychologie de la violence}\theme{société!organisation sociale}\theme{paix!pacifisme}\theme{idéologie}}

%%%%%%%%%%%%%%%%%%%%%%%%%%%%%Nv-Sujet%%%%%%%%%%%%%%%%%%%%%%%%%%%%%%%%%%%%%%%%%%
\sujetn{%
	nomauteur=Popper,
	prenomauteur=Karl,
	initialeprenom =K,
	particule=,
	titre = {La société ouverte et ses ennemis},
	semestre = \HQ,
	lieu ={Amérique du Nord},
	annee ={2023},
	type =,
	jour =1,
	datepub =1962,
	ajoutref ={traduction Jacqueline \textsc{Bernard} et Philippe \textsc{Monod}},
	genre =\essaip,
	corrige =Oui,
	liencorrige=,
	typequn = \phil,
	typeqdeux = \litt,
	qun = {Pensez-vous que Popper démontre ici que la société contemporaine est déshumanisante~?},
	qdeux ={La littérature peut-elle nous aider à bien vivre en société~?},
	themeprogrammequn=HL,
	themeprogrammeqdeux=OdH}
{Imaginons, au prix d'une
certaine exagération, une société où les hommes ne se rencontrent
jamais face à face, où les affaires sont traitées par des individus
isolés communiquant entre eux par lettres ou par télégrammes, se
déplaçant en voiture fermée et se reproduisant par insémination
artificielle~: pareille société serait totalement abstraite et
dépersonnalisée. Or la société moderne lui ressemble déjà sur bien
des points. Dans une ville, les piétons se croisent mais s'ignorent,
les membres d'un syndicat portent une carte et paient une cotisation
mais peuvent ne jamais se connaître. Beaucoup d'individus ont peu ou
pas de contacts humains et vivent dans l'anonymat et l'isolement. Ils
y sont malheureux, car, si la société tend à devenir abstraite, le
tissu biologique de l'homme n'a guère changé, et il a des besoins
sociaux que celle-ci est incapable de satisfaire.

Certes, il n'y aura jamais de société complètement abstraite, pas
plus, d'ailleurs, qu'on n'en connaîtra entièrement ou
d'essentiellement rationnelle. Les hommes continuent à former des
groupes concrets et à maintenir entre eux divers rapports sociaux.
Mais, à l'exception de ce qui se passe dans certaines cellules
familiales, la plupart de ces groupes ne s'ouvrent sur aucune vie
commune et beaucoup ne jouent aucun rôle sur le plan social.

Il faut cependant être juste. L'évolution de la société présente aussi
des aspects positifs. Les rapports personnels, n'étant plus
déterminés par le hasard de la naissance, peuvent s'épanouir
librement en une nouvelle forme d'individualisme. De même, les liens
spirituels peuvent jouer un rôle d'autant plus important que les
liens physiques ou biologiques sont affaiblis.
\theme{déshumanisation}\theme{technique}\theme{technique!et sciences}\theme{eugénisme}\theme{naturel (et artificiel)}\theme{dépersonnalisation}\theme{société!organisation sociale}\theme{biogénétique}\theme{biologie}\theme{société!société abstraite}\theme{famille}\theme{individu}\theme{individualisme}}

%%%%%%%%%%%%%%%%%%%%%%%%%%%%%Nv-Sujet%%%%%%%%%%%%%%%%%%%%%%%%%%%%%%%%%%%%%%%%%%
\sujetn{%
	nomauteur=Barbusse,
	prenomauteur=Henri,
	initialeprenom =H,
	particule=,
	titre = {Le Feu},
	semestre = \HQ,
	lieu ={Amérique du Nord},
	annee =2023,
	type =,
	jour =2,
	datepub =1916,
	ajoutref =,
	genre =\romanl,
	corrige =,
	typequn = \litt,
	typeqdeux = \phil,
	qun = {Quelle vision paradoxale de la guerre se dégage de cet extrait~?},
	qdeux ={Peut-on donner une description objective de la guerre~?},
	themeprogrammequn=HV,
	themeprogrammeqdeux=HV}
{\chapeau{Dans ce roman qui se
déroule pendant la Première guerre mondiale, l'auteur nous libre sa
propre expérience du front.}

On se remet en marche, parsemés sur la route maintenant grisâtre, très
lentement, très pesamment, avec des geignements\footnote{Geignements
: petits cris plaintifs.} et de sourdes malédictions que l'effort
étrangle dans les gorges. Au bout de cent mètres, les deux hommes
formant équipe échangent leurs fardeaux, de sorte qu'au bout de deux
cents mètres, malgré la bise aigre\footnote{Bise aigre~: vent
glacial.} et blanchissante du petit matin, tout le monde, sauf les
gradés, ruisselle de sueur.

Tout à coup une étoile intense s'épanouit là-bas, vers les lieux
vagues où nous allons~: une fusée. Elle éclaire toute une portion du
firmament de son halo\footnote{Halo~: cercle lumineux.} laiteux, en
effaçant les constellations, et elle descend gracieusement, avec des
airs de fée.

Une rapide lumière en face de nous, là-bas~; un éclair, une
détonation.

C'est un obus~!

Au reflet horizontal que l'explosion a instantanément répandu dans le
bas du ciel, on voit nettement que, devant nous, à un kilomètre
peut-être, se profile, de l'est à l'ouest, une crête.

Cette crête est à nous dans toute la partie visible d'ici, jusqu'au
sommet, que nos troupes occupent. Sur l'autre versant, à cent mètres
de notre première ligne, est la première ligne allemande.

L'obus est tombé sur le sommet, dans nos lignes. Ce sont eux qui
tirent.

Un autre obus. Un autre, un autre, plantent, vers le haut de la
colline, des arbres de lumière violacée dont chacun illumine
sourdement tout l'horizon.

Et bientôt, il y a un scintillement d'étoiles éclatantes et une forêt
subite de panaches phosphorescents sur la colline~: un mirage de
féerie bleu et blanc se suspend légèrement à nos yeux dans le gouffre
entier de la nuit.

Ceux d'entre nous qui consacrent toutes les forces arc-boutées de
leurs bras et de leurs jambes à empêcher leurs vaseux\footnote{Vaseux
: recouverts de vase, de boue.} fardeaux trop lourds de leur glisser
du dos et à s'empêcher eux-mêmes de glisser par terre, ne voient rien
et ne disent rien. Les autres, tout en frissonnant de froid, en
grelottant, en reniflant, en s'épongeant le nez avec des mouchoirs
mouillés\footnote{Mouchoirs mouillés~: à cause du manque de masques à
gaz, les soldats français avaient reçu l'ordre de se protéger le nez
et la bouche avec des mouchoirs humides.} qui pendent de l'aile, en
maudissant les obstacles de la route en lambeaux, regardent et
commentent.

—~C'est comme si tu vois un feu d'artifice, disent-ils.

Complétant l'illusion de grand décor d'opéra féerique et sinistre
devant lequel rampe, grouille et clapote notre troupe basse, toute
noire, voici une étoile rouge, une verte~; une gerbe rouge, beaucoup
plus lente.

On ne peut s'empêcher, dans nos rangs, de murmurer avec un confus
accent d'admiration populaire, pendant que la moitié disponible des
paires d'yeux regardent~:

—~Oh~! une rouge~!… Oh~! une verte~!…
\theme{guerre}\theme{armée}\theme{bataille}\theme{marche (militaire)}\theme{esthétique!de la violence}\theme{combat}\theme{armes}}

%%%%%%%%%%%%%%%%%%%%%%%%%%%%%Nv-Sujet%%%%%%%%%%%%%%%%%%%%%%%%%%%%%%%%%%%%%%%%%%
\sujetn{%
	nomauteur=Barjavel,
	prenomauteur=René,
	initialeprenom =R,
	particule=,
	titre = {Ravage},
	semestre = \HQ,
	lieu =Centres étrangers,
	annee =2023,
	type =,
	jour =1,
	datepub =1943,
	ajoutref =,
	genre =\romanl,
	corrige =Oui,
%	liencorrige={cea_2023_1},
	typequn = \litt,
	typeqdeux = \phil,
	qun = {Quelle image de la science cet extrait propose-t-il~?},
	qdeux ={Doit-on attendre de la science qu'elle réponde à toutes nos questions~?},
	themeprogrammequn=HL,
	themeprogrammeqdeux=HLCC}
{\chapeau{Dans ce roman paru en 1943, l'auteur imagine une société dans laquelle
l'électricité disparaît, ce qui plonge le pays dans le chaos.}

La foule se serra autour de la voiture à bras. Elle se sentait
rassurée par la présence de cet homme qui connaissait les secrets de
la nature. Chacun avait l'impression de se retrouver près de la
Science elle-même, la Science qui explique tout et peut tout. Un
monsieur maigre s'empara d'un seau en fer que portait une ménagère,
le posa à terre, renversé, grimpa dessus, et, d'une voix de coq
enroué, parla~:

—~Messieurs, mesdames, citoyens…

—~Hou… hou…, répondit la foule.

—~Je ne veux pas vous faire de discours, je me propose seulement de
demander en votre nom à l'éminent savant qui se retrouve en ce moment
parmi nous de nous donner des éclaircissements sur le phénomène qui
vient de bouleverser notre vie. Je…

—~Vive Portin~! La parole à Portin. Portin~! Portin~! Portin~!

Le savant tremblait d'émotion dans son fauteuil et faisait avec la
main des gestes de dénégation. Alors un gigantesque ouvrier fendit
les groupes et parvint jusqu'à la charrette. C'était un
métallurgiste, un ancien du métier, à la peau recuite, un vieux
compagnon qui avait résisté à trente ans d'usine. Sa main droite,
avec laquelle, à l'atelier, il donnait toutes les deux secondes le
même coup de marteau sur des rivets toujours pareils, restait fermée
autour d'un manche imaginaire.

—~Écoutez, m'sieur Portin, nous, on est là, on sait pas, et on veut
savoir. Vous, vous savez, la Science… Il faut nous dire. Qu'est-ce
qui se passe~? Quand est-ce que ça va finir~?

Le vieillard, péniblement, se leva de son fauteuil. Il tremblait.

—~Mes bons amis… dit-il

Sa voix aigrelette ne portait pas à plus de dix mètres.

—~Mes bons amis, je ne peux rien vous dire, je ne sais rien. On n'a
jamais vu ça.

Notre science est une science expérimentale. Or, le phénomène qui
vient de se produire ne correspond à rien de ce que nous savons.
C'est en violant toutes les lois de la Nature et la logique que
l'électricité a disparu. Et, l'électricité morte, il est encore plus
invraisemblable que nous soyons vivants. Tout cela est fou. C'est un
cauchemar antiscientifique, antirationnel. Toutes nos théories,
toutes nos lois sont renversées. Voir cela au terme de ma vie de
savant…

Il se laissa retomber lourdement dans son fauteuil. Les premiers rangs
de la foule virent de grosses larmes couler de ses yeux dans sa
moustache blanche. Mais les gens qui se trouvaient plus loin,
inquiets, curieux, voulurent aussi entendre. Les grands se haussaient
sur la pointe des pieds, les petits se cramponnaient aux grands. Des
gamins grimpaient aux fûts des lampadaires. On se passait de rang à
rang des fragments de phrase~:

—~Il a dit que l'électricité était morte.

—~Ma pauvre, il a dit qu'il y comprenait rien.

—~Il a dit que c'était la guerre.

—~Il a dit qu'il allait tout arranger.

La multitude voulut en savoir davantage. De partout à la fois elle
poussa vers le centre. Dix mille poitrines firent pression. La foule
ne fut plus qu'une masse compacte, un seul muscle contracté. Il y eut
des remous, des tourbillons, des vêtements arrachés, des côtes
fracturées, des caleçons souillés. La voiture de M. Paul Portin fit
trois tours sur elle-même, craqua et disparut. Le vieux savant se
trouva projeté en l'air et retomba sur des épaules. Il y flotta
quelques instants, puis sombra.
\theme{sciences}\theme{technologie}\theme{électricité}\theme{dystopie}\theme{nature!lois de la nature}\theme{société}\theme{scientifique}\theme{technique!et sciences}\theme{technique}}

%%%%%%%%%%%%%%%%%%%%%%%%%%%%Nv-Sujet%%%%%%%%%%%%%%%%%%%%%%%%%%%%%%%%%%%%%%%%%%

\sujetn{%
	nomauteur=Ponthus,
	prenomauteur=Joseph,
	initialeprenom =J,
	particule=,
	titre = {À la ligne, Feuillets d'usine},
	semestre = \HQ,
	lieu =Métropole,
	annee =2023,
	type =Remplacement,
	jour =2,
	datepub =2019,
	ajoutref =,
	genre =\poesie,
	corrige =Oui,
%	liencorrige={metrem_2023_2},
	typequn = \litt,
	typeqdeux = \phil,
	qun = {«~Je suis là sans y être.~», dit le poète. En quoi cette phrase rend-elle compte du travail à l'usine tel qu'il est décrit dans ce texte~?},
	qdeux ={Le travail menace-t-il notre dignité~?},
	themeprogrammequn=HL,
	themeprogrammeqdeux=HL}
{\chapeau{À la ligne est un roman qui présente l'allure d'un poème. Il évoque l'histoire d'un ouvrier intérimaire qui, pour se rapprocher de sa femme, trouve un emploi dans les conserveries de poissons.}

\setlength{\parskip}{0pt}

Le matin c'est la nuit

L'après-midi c'est la nuit

La nuit c'est encore pire

\medskip

Dès qu'on rentre dans l'usine c'est la nuit

Les néons

L'absence de fenêtres dans tous nos immenses cubes d'ateliers

Une nuit qui va durer nos huit heures de travail minimum

\medskip

On sort du sommeil encore marqué de rêves d'usine

Pour replonger dans une autre nuit

Artificielle froide et éclairée de néons

\medskip

Dès lors

C'est comme si

On continuait sa nuit

\medskip

Entre la nuit de la maison et celle de l'embauche

Le réveil

Deux heures de transition

Yeux embrumés et cafés serrés

Ce serait donc ça le matin

\medskip

Tous les matins du monde

\medskip

Au vestiaire avant l'embauche

Encore cinq minutes avant de replonger dans la nuit

J'admire sincèrement toutes les ouvrières et tous les ouvriers qui ont pris 
leur douche et se sont parfumés

Moi je ne peux pas

La douche c'est le soir

Enfin en rentrant du boulot

Certains se pomponnent vraiment

Doivent s'affairer devant leur miroir

\medskip

Les miroirs à l'usine

Il y en a aux intersections et aux virages à angle droit dans les couloirs 
sans fin de sorte que l'on voie les transpalettes arriver et ainsi ne pas se 
faire rouler dessus

On ne se regarde pas dedans

De toute façon on sait à quoi l'on ressemble

Une tenue blanche maculée de sang comme celle de tout le monde

Des corps fatigués

Des yeux qui continuent

\medskip

Pas moi

L'usine c'est une tenue que je garde une semaine et qui se cradifie et pue de plus en plus au fil des jours

Machine le vendredi ou le samedi

\medskip

J'imagine que pour elles et pour eux

Ce doit être affaire de dignité ou d'une certaine forme de noblesse

Que d'arriver propre et sentant bon à l'abattoir

Malgré ce qu'il y aura à faire

\medskip

Le matin

Entre mes deux nuits

Je suis là sans y être

Comme si

J'étais en transition

La vraie vie sera à la débauche\footnote{la débauche~:\\sens 1~: usage excessif ou déréglé de tous les plaisirs des sens, notamment ceux de l'amour et de la table\\sens 2~: par opposition à l'embauche, moment où l'on cesse le travail --- débauchage}
\setlength{\parskip}{6pt}
\theme{ouvriers!condition ouvrière}\theme{aliénation}\theme{déshumanisation}\theme{naturel (et artificiel)}\theme{artificiel|see{naturel (et artificiel)}}\theme{usine}\theme{corps}\theme{dignité humaine}\theme{machines}\theme{nuit}}

%%%%%%%%%%%%%%%%%%%%%%%%%%%%Nv-Sujet%%%%%%%%%%%%%%%%%%%%%%%%%%%%%%%%%%%%%%%%%%

\sujetn{%
	nomauteur=Ponthus,
	prenomauteur=Joseph,
	initialeprenom =J,
	particule=,
	titre = {À la ligne, Feuillets d'usine},
	semestre = \HQ,
	lieu =Nouvelle-Calédonie,
	annee =2023,
	type =,
	jour =1,
	datepub =2019,
	ajoutref =,
	genre =\poesie,
	corrige =,
	typequn = \litt,
	typeqdeux = \phil,
	qun = {Comment s'exprime, dans ce texte, la violence de l'usine sur «~toutes les ouvrières et tous les ouvriers~»~?},
	qdeux ={Peut-on rester humain dans des conditions inhumaines~?},
	themeprogrammequn=HV,
	themeprogrammeqdeux=HVL}
{\chapeau{À la ligne est un roman qui présente l'allure d'un poème. Il évoque l'histoire d'un ouvrier intérimaire qui, pour se rapprocher de sa femme, trouve un emploi dans les conserveries de poissons.}

\setlength{\parskip}{0pt}

Le matin c'est la nuit

L'après-midi c'est la nuit

La nuit c'est encore pire

\medskip

Dès qu'on rentre dans l'usine c'est la nuit

Les néons

L'absence de fenêtres dans tous nos immenses cubes d'ateliers

Une nuit qui va durer nos huit heures de travail minimum

\medskip

On sort du sommeil encore marqué de rêves d'usine

Pour replonger dans une autre nuit

Artificielle froide et éclairée de néons

\medskip

Dès lors

C'est comme si

On continuait sa nuit

\medskip

Entre la nuit de la maison et celle de l'embauche

Le réveil

Deux heures de transition

Yeux embrumés et cafés serrés

Ce serait donc ça le matin

Tous les matins du monde

\medskip

Au vestiaire avant l'embauche\footnote{\emph{Embauche}~: l'embauche désigne le moment où l'on commence le travail.}

Encore cinq minutes avant de replonger dans la nuit

J'admire sincèrement toutes les ouvrières et tous les ouvriers qui ont pris 
leur douche et se sont parfumés

Moi je ne peux pas

La douche c'est le soir

Enfin en rentrant du boulot

Certains se pomponnent vraiment

Doivent s'affairer devant leur miroir

\medskip

Les miroirs à l'usine

Il y en a aux intersections et aux virages à angle droit dans les couloirs 
sans fin de sorte que l'on voie les transpalettes arriver et ainsi ne pas se 
faire rouler dessus

On ne se regarde pas dedans

De toute façon on sait à quoi l'on ressemble

Une tenue blanche maculée de sang comme celle de tout le monde

Des corps fatigués

Des yeux qui continuent

\medskip

Pas moi

L'usine c'est une tenue que je garde une semaine et qui se cradifie et pue de plus en plus au fil des jours

Machine le vendredi ou le samedi

\medskip

J'imagine que pour elles et pour eux

Ce doit être affaire de dignité ou d'une certaine forme de noblesse

Que d'arriver propre et sentant bon à l'abattoir

Malgré ce qu'il y aura à faire

\medskip

Le matin

Entre mes deux nuits

Je suis là sans y être

Comme si

J'étais en transition

La vraie vie sera à la débauche\footnote{\emph{Débauche}~: le mot s'emploie ici comme l'antonyme d'\emph{embauche}. La débauche est le
moment où l'on quitte le travail.
\setlength{\parskip}{6pt}}
\theme{ouvriers!condition ouvrière}\theme{aliénation}\theme{déshumanisation}\theme{naturel (et artificiel)}\theme{usine}\theme{corps}\theme{dignité humaine}\theme{machines}\theme{nuit}}

%%%%%%%%%%%%%%%%%%%%%%%%%%%%Nv-Sujet%%%%%%%%%%%%%%%%%%%%%%%%%%%%%%%%%%%%%%%%%%

\sujetn{%
	nomauteur=Russell,
	prenomauteur=Bertrand,
	initialeprenom =B,
	particule=,
	titre = {Essais sceptiques},
	semestre = \HQ,
	lieu =Polynésie,
	annee =2023,
	type =,
	jour =2,
	datepub =1928,
	ajoutref ={«~Introduction~: de la valeur du scepticisme~»},
	genre =\essaip,
	corrige =,
	typequn = \phil,
	typeqdeux = \litt,
	qun = {Selon Russell, en quoi la guerre est-elle une folie~?},
	qdeux ={La littérature et les arts éclairent-ils les mécanismes de la violence~?},
	themeprogrammequn=HV,
	themeprogrammeqdeux=HV}
{Les événements humains
naissent des passions, qui engendrent les systèmes de mythes qui les
accompagnent. Les psychanalystes ont étudié les manifestations
individuelles de ce processus chez les fous, reconnus ou non reconnus
comme tels. Un homme qui a souffert quelque humiliation invente une
théorie selon laquelle il est roi d'Angleterre et il trouve toutes
sortes d'explications ingénieuses pour justifier le fait qu'il n'est
pas traité avec tous les égards dus à sa haute situation. Dans ce cas
particulier, son illusion ne provoque pas de sympathie de la part de
ses voisins, et c'est pourquoi ils l'enferment. Mais si, au lieu
d'affirmer sa propre grandeur, il affirme la grandeur de sa nation ou
de sa classe ou de sa foi, il gagne des armées d'adhérents et devient
un chef politique ou religieux, même si, pour un observateur
impartial, son opinion semble aussi absurde que celles qu'on trouve
dans les asiles de fous. Ainsi naissent des folies collectives qui
suivent des lois très analogues à celles de la folie individuelle.
Tout le monde sait qu'il est dangereux de discuter avec un fou qui se
croit roi de l'Angleterre~; mais comme il est isolé, on peut avoir
raison de lui. Quand une nation tout entière participe à une erreur,
sa colère est du genre de celle d'un individu fou quand on discute
ses prétentions~; mais il ne faut pas moins d'une guerre pour la
forcer à se soumettre à la raison.
\theme{guerre}\theme{guerre!causes de la guerre}\theme{folie}\theme{classes sociales}\theme{religion}\theme{patriotisme}\theme{société}}

%%%%%%%%%%%%%%%%%%%%%%%%%%%%Nv-Sujet%%%%%%%%%%%%%%%%%%%%%%%%%%%%%%%%%%%%%%%%%%

\sujetn{%
	nomauteur=Nietzsche,
	prenomauteur=Friedrich,
	initialeprenom =F,
	particule=,
	titre = {Le Gai Savoir},
	semestre = \HQ,
	lieu ={Amérique du Nord},
	annee =2024,
	type =,
	jour =2,
	datepub =1882,
	ajoutref ={(trad. Alexandre \textsc{Vialatte})},
	genre =\essaip,
	corrige =Oui,
%	liencorrige={amnord_2024_2},
	typequn = \phil,
	typeqdeux = \litt,
	qun = {Comment l'auteur défend-il dans ce texte le «~nouveau~», contre l'opinion commune~?},
	qdeux ={Selon vous, le «~goût du neuf, du risqué, de l'inessayé~» suffit-il à définir le rôle de l'écrivain ou de l'artiste~?},
	themeprogrammequn=CC,
	themeprogrammeqdeux=CC}
{Ce sont les esprits forts et les esprits malins, les plus forts et les plus malins qui ont fait faire jusqu'ici le plus de progrès à l'humanité~: ils ont rallumé constamment les passions qui allaient s'endormir --- toute société policée les endort ---, ils ont réveillé constamment l'esprit de comparaison et de contradiction, le goût du neuf, du risque, de l'inessayé~; ils ont obligé l'homme à opposer sans cesse les opinions aux opinions, les idéaux aux idéals. Par les armes le plus souvent, en renversant les bornes-frontières, en violant les piétés, mais aussi en fondant de nouvelles religions, en créant de nouvelles morales~! Cette «~méchanceté~» qu'on retrouve dans tout professeur de nouveau, dans tout prédicateur de choses neuves, c'est la même «~méchanceté~» qui discrédite le conquérant, bien qu'elle s'exprime plus subtilement et ne mobilise pas immédiatement le muscle~; --- ce qui fait d'ailleurs qu'elle discrédite moins fort~! --- Le neuf, de toute façon, c'est le mal, puisque c'est ce qui veut conquérir, renverser les bornes-frontières, abattre les anciennes pitiés~; seul l'ancien est le bien~! Les hommes de bien, à toute époque, sont ceux qui plantent profondément les vieilles idées pour leur faire porter du fruit, ce sont les cultivateurs de l'esprit. Mais tout terrain finit par s'épuiser, il faut toujours que la charrue du mal y revienne.
\theme{armes}\theme{art}\theme{nouveauté}\theme{progrès}\theme{passion}\theme{idéal}\theme{art!religion}\theme{morale}\theme{ancien}\theme{culture}\theme{pitié}}

%%%%%%%%%%%%%%%%%%%%%%%%%%%%Nv-Sujet%%%%%%%%%%%%%%%%%%%%%%%%%%%%%%%%%%%%%%%%%%

\sujetn{%
	nomauteur=Aragon,
	prenomauteur=Louis,
	initialeprenom =L,
	particule=,
	titre = {Le Roman inachevé},
	semestre = \HQ,
	lieu =Métropole,
	annee =2024,
	type =,
	jour =1,
	datepub =1956,
	ajoutref =,
	genre =\poesie,
	corrige =Oui,
%	liencorrige={met_2024_1},
	typequn = \litt,
	typeqdeux = \phil,
	qun = {Quelles violences ce poème dit-il~?},
	qdeux ={Peut-on perdre son humanité~?},
	themeprogrammequn=HV,
	themeprogrammeqdeux=HVL}
{\setlength{\parskip}{0pt}
\begin{center}
\textbf{Classe 17\footnote{Une classe est l'ensemble des hommes aptes au service militaire nés la même année. La classe 17 désigne donc les hommes de vingt ans qui ont été recensés en 1917.}}
\end{center}

\bigskip

Je me souviens C'était je crois tout près de Saint-Michel-en-Grève

Mais peut-être après tout que je confonds la vie avec le rêve

\bigskip

Je ne savais pas qu'on pût traiter ainsi des êtres humains

\bigskip

Je me souviens Il était venu des gens de tous les villages

On en voyait arriver au loin par les sables de la plage

\bigskip

Il y avait des groupes de paysans sur tous les chemins

\bigskip

Des villégiateurs\footnote{«~villégiateurs~»~: ceux qui séjournent dans un 
lieu de vacances ou de loisirs.\vfill{~}} avec leur marmaille de toile blanche

Des bourgeois de Lannion qui poussaient jusque-là le dimanche

\bigskip

Un monde au bord de la mer avec des chapeaux noirs et des gants

\bigskip

De plus en plus le temps gris de l'été tournait à l'étouffoir

Avec tout ça les coiffes conféraient aux prés un air de foire

\bigskip

Le ciel portait un manège d'oiseaux criards et fatigants

\bigskip

Tout à coup les pêcheurs abandonnent leurs filets dans les roches

Les jambes sont pleines d'enfants qui courent et crient qu'ils approchent

\bigskip

Et les voilà chargés de poussière et d'humiliation

\bigskip

Troupeau confus les boutons arrachés aux capotes de terre

Sans armes sans ceinturons enroués à force de se taire

\bigskip

Le pas rompu le visage étrangement sans expression

\bigskip

Couleur des murs longés les yeux gris une barbe de trois jours

Blonde ou rousse et le regard égaré des fauves et des sourds

\bigskip

Les voilà comme un cheminement maudit dans des champs pierreux

\vfill{~}\pagebreak

Plus grands que nature à côté des fusils Gras\footnote{Le fusil Gras est employé par l'armée française jusqu'en 1887. Ensuite, cette arme n'est plus utilisée que de manière marginale, en particulier par les pompiers, les gardes républicains, etc.} qui les escortent

À travers ce pays où pour eux les maisons n'ont pas de portes

\bigskip

Ce pays qui n'a que des bornes kilométriques pour eux

\bigskip

Ils ont la tête qui retentit toujours des tirs de barrage

Et trop de poux qu'on leur permette de dormir dans le fourrage

\bigskip

C'est près de Reims qu'on les a pris comme des mouches dans la craie

\bigskip

Cette terre a le crâne dur On a bien du mal à s'y faire

Elle a gardé morts et vivants à son abri tout un hiver

\bigskip

Mais un beau matin de printemps en a livré tous les secrets

\bigskip

Depuis ce jour leur long malheur s'étire comme une couleuvre

Ils ne sont que des prisonniers que l'on achemine à pied d'œuvre

\bigskip

Ils ont marché marché marché comme ils vivaient dans les tranchées

\bigskip

Ils ont marché marché marché jusqu'au-delà de la fatigue

Les pieds et la mémoire en sang rêvant la Saxe ou le Schlesvig\footnote{«~la 
Saxe ou le Schlesvig~»~: régions d'Allemagne}

\bigskip

Et sans savoir où ils allaient ils ont marché marché marché

\bigskip

Après tout les voilà contents d'être sortis de la bataille

Des fermiers tâtent leurs mollets pour voir si c'est du bon bétail

\bigskip

On a des morts dans la commune on les remplacera comment

\bigskip

Celui-là tenez le rouquin nous servirait pour les cultures

Est-ce qu'ils sont très exigeants sur la question nourriture

\bigskip

Avec tous ceux qui sont partis on prendrait bien des Allemands

\setlength{\parskip}{6pt}
\theme{esthétique!de la violence}\theme{armes}\theme{guerre}\theme{combat}\theme{déshumanisation}\theme{guerre!civils}\theme{enfant}\theme{marche (militaire)}\theme{Poilus}\theme{guerre!Première Guerre mondiale}\theme{vie (et mort)}\theme{mort|see{vie (et mort)}}}

%%%%%%%%%%%%%%%%%%%%%%%%%%%Nv-Sujet%%%%%%%%%%%%%%%%%%%%%%%%%%%%%%%%%%%%%%%%%

\sujetn{%
	nomauteur=Aragon,
	prenomauteur=Louis,
	initialeprenom =L,
	particule=,
	titre = {Le Roman inachevé},
	semestre = \HQ,
	lieu =Polynésie,
	annee =2024,
	type =,
	jour =1,
	datepub =1956,
	ajoutref =,
	genre =\poesie,
	corrige =,
	typequn = \litt,
	typeqdeux = \phil,
	qun = {Quelle image le poète donne-t-il de ses camarades disparus~?},
	qdeux ={La violence efface-t-elle toute trace d'humanité~?},
	themeprogrammequn=HV,
	themeprogrammeqdeux=HV}
	%%%Texte%%%
	{\chapeau{Durant les deux guerres mondiales, le poète Louis Aragon fut mobilisé comme médecin auxiliaire sur la ligne de front. Il évoque ici la figure de ses camarades morts pour la patrie.}
\setlength{\parskip}{0pt}
	\begin{center}
	[…]
	\end{center}
		
Tu n'en reviendras pas toi qui courais les filles

Jeune homme dont j'ai vu battre le cœur à nu

Quand j'ai déchiré ta chemise et toi non plus

Tu n'en reviendras pas vieux joueur de manille\footnote{La manille est un jeu de cartes d'origine espagnole.}

\bigskip

Qu'un obus a coupé par le travers en deux

Pour une fois qu'il avait un jeu du tonnerre

Et toi le tatoué l'ancien Légionnaire

Tu survivras longtemps sans visage sans yeux

\bigskip

Roule au loin roule train des dernières lueurs

Les soldats assoupis que ta danse secoue

Laissent pencher leur front et fléchissent le cou

Cela sent le tabac la laine et la sueur

\bigskip

Comment vous regarder sans voir vos destinées

Fiancés de la terre et promis des douleurs

La veilleuse\footnote{Une veilleuse est une petite lampe qu'on laisse allumée durant la nuit.} vous fait de la couleur des pleurs

Vous bougez vaguement vos jambes condamnées

\bigskip

Vous étirez vos bras vous retrouvez le jour

Arrêt brusque et quelqu'un crie Au jus là-dedans\footnote{Argot militaire~: par cette expression, le caporal appelle les soldats à venir boire le café.}

Vous bâillez vous avez une bouche et des dents

Et le caporal chante Au pont de Minaucourt\footnote{Au pont de Minaucourt est une chanson militaire qui évoque une violente bataille de la Première Guerre mondiale.}

\bigskip

Déjà la pierre pense où votre nom s'inscrit

Déjà vous n'êtes plus qu'un mot d'or sur nos places

Déjà le souvenir de vos amours s'efface

Déjà vous n'êtes plus que pour avoir péri
\setlength{\parskip}{6pt}

\theme{vie (et mort)}\theme{guerre}\theme{blessure}\theme{guerre!Première Guerre mondiale}\theme{guerre!Seconde Guerre mondiale}\theme{bataille}\theme{mort}\theme{victime}\theme{soldats}\theme{souffrance}\theme{armes}\theme{gueules cassées}\theme{corps}\theme{chanson}\theme{mémoire}}

%%%%%%%%%%%%%%%%%%%%%%%%%%%Nv-Sujet%%%%%%%%%%%%%%%%%%%%%%%%%%%%%%%%%%%%%%%%%

\sujetn{%
	nomauteur=Delbo,
	prenomauteur=Charlotte,
	initialeprenom =C,
	particule=,
	titre = {Mesure de nos jours},
	semestre = \HQ,
	lieu =Amérique du Nord,
	annee =2025,
	type =,
	jour =2,
	datepub =1971,
	ajoutref =,
	genre =poesie,
	corrige =,
%	liencorrige=,
	typequn = \litt,
	typeqdeux = \phil,
	qun = {Comment l'autrice exprime-t-elle ici l'impossible retour à la vie après avoir subi la violence de l'histoire~?},
	qdeux ={Suffit-il de ne plus subir la violence pour y échapper~?},
	themeprogrammequn=HV,
	themeprogrammeqdeux=HV}
	%%%Texte%%%
	{\setlength{\parskip}{0pt}
	
Rentrer du camp rentrer dans le rang
	
après l'histoire

le tous les jours

après le maquis\footnote{Maquis~: lieu où se réunissaient les résistants à l'occupation allemande pendant la Seconde Guerre mondiale.}

le traintrain de la vie.

Nous disions

que la vie sera belle quand elle sera libre

que la vie sera ardente quand nous serons libres

tout sera simple

transparent

tout nous sera rendu avec la liberté

la beauté l'amour l'amitié

tout

la liberté

c'est tout

il n'y aura qu'à vivre

quoi de plus simple

de plus facile

à celui qui sait souffrir

à celui qui sait mourir~?

Rentrer

Qui de nous osait penser plus loin~?

Rentrer

c'était déjà demander l'impossible

c'était tout demander

oserait-on demander davantage~?

Rentrer

tout nous serait rendu.


Revenir, ce n'est pas tout

c'est revenir pour se remettre à vivre

à vivre le tous les jours

à travailler et à faire des dettes

à faire des économies pour payer ses dettes

à vendre du savon

parce qu'on ne sait pas faire autre chose

à retourner au bureau

parce qu'on ne sait pas faire autre chose

dans la vie de tous les jours

à chercher un logement

parce qu'on ne peut pas vivre autrement

à être à l'heure

parce qu'au travail il faut être à l'heure.

…

De quoi vous plaignez-vous

la vie c'est la vie

de quoi rêviez-vous dans votre là-bas~?

De manger à votre faim

de dormir à votre sommeil

d'aimer à votre amour

À manger à dormir à aimer

vous l'avez

depuis que vous êtes rentrés.

L'histoire

c'est fini

soyez heureux comme tout le monde

l'histoire

c'est un moment

maintenant

c'est la vie.

Et pourquoi donc vouliez-vous revenir~?

\medskip

Sortir de l'histoire

pour entrer dans la vie

essayez donc vous autres et vous verrez.
	\setlength{\parskip}{6pt}
\theme{camp de concentration}\theme{Résistants}\theme{liberté}\theme{vie (et mort)!sens de la vie}\theme{quotidien}\theme{déshumanisation}\theme{histoire}\theme{bonheur}}

%%%%%%%%%%%%%%%%%%%%%%%%%%%Nv-Sujet%%%%%%%%%%%%%%%%%%%%%%%%%%%%%%%%%%%%%%%%%

\sujetn{%
	nomauteur=Emerson,
	prenomauteur=Ralph Waldo,
	initialeprenom =R. W,
	particule=,
	titre = {L'intellectuel américain},
	semestre = \HQ,
	lieu =Centres étrangers,
	annee =2025,
	type =,
	jour =2,
	datepub =1837,
	ajoutref =,
	genre =\essaip,
	corrige =oui,
%	liencorrige={cea_2025_2},
	typequn = \phil,
	typeqdeux = \litt,
	qun = {Quelle distinction Emerson propose-t-il entre l'«~Homme Pensant~» et le «~rat de bibliothèque~»~?},
	qdeux ={Faut-il rompre avec la tradition littéraire pour écrire une œuvre~?},
	themeprogrammequn=ETACC,
	themeprogrammeqdeux=CC}
	%%%Texte%%%
	{Chaque époque, on l'aura compris, doit écrire ses propres livres~; mieux, chaque génération doit écrire pour celle qui la suivra. Les livres d'une période plus ancienne n'entrent pas dans ce cadre.
	
Cependant, c'est ici que se produit une grave erreur. Le caractère sacré qui s'attache à l'acte de création, à l'acte de pensée, est transféré sur le produit de ces actes. Le barde\footnote{Barde~: poète celtique qui célébrait en chantant les exploits des héros.} était considéré comme un homme divin~: donc, son chant est également divin. L'écrivain était vu comme un esprit juste et sage~: donc, il s'ensuit que son livre est parfait~; tout comme l'amour que l'on ressent pour le héros est ensuite corrompu en vénération idolâtre de la statue. Immédiatement, le livre devient nocif et le guide se fait tyran. L'esprit léthargique\footnote{Léthargique~: ensommeillé, engourdi.} et perverti de la multitude, si lent à s'ouvrir aux incursions\footnote{Incursion~: surgissement momentané.} de la Raison, une fois qu'il est ainsi ouvert, une fois qu'il a reçu le livre, se repose dessus et se répand en protestations si ce livre est critiqué. On bâtit des universités sur ce livre. Des penseurs, mais non l'«~Homme Pensant~» écrivent d'autres livres sur celui-là~; ces livres sont certes écrits par des hommes de talent, mais ces hommes partent sur de mauvaises bases, ils réfléchissent à partir de dogmes acceptés et non à partir de leurs propres conceptions des grands principes. C'est ainsi que grandissent dans les bibliothèques des jeunes gens humbles et mous, qui croient qu'il est de leur devoir d'accepter les vues que Cicéron\footnote{Cicéron~: philosophe du 1er siècle avant J.-C.}, Locke\footnote{Locke~: philosophe anglais (1632-1704).} ou Bacon\footnote{Bacon~: philosophe anglais (1561-1626).} ont exposées~; ils oublient que Cicéron, Locke et Bacon n'étaient que des jeunes gens eux-mêmes enfermés dans des bibliothèques, quand ils ont écrit leurs livres.

Et alors, au lieu de l'«~Homme Pensant~», nous aurons droit au rat de bibliothèque.
\theme{création}\theme{lecture}\theme{génie}\theme{écrivain}\theme{écriture}\theme{nouveauté}\theme{tradition}\theme{critique (littéraire)}\theme{influence}\theme{génération}\theme{penseur}
}

%%%%%%%%%%%%%%%%%%%%%%%%%%%Nv-Sujet%%%%%%%%%%%%%%%%%%%%%%%%%%%%%%%%%%%%%%%%%

\sujetn{%
	nomauteur=Szymborska,
	prenomauteur=Wislawa,
	initialeprenom =W,
	particule=,
	titre = {Vue avec grain de sable},
	semestre = \HQ,
	lieu =Métropole,
	annee =2025,
	type =,
	jour =2,
	datepub ={1996},
	ajoutref =\linebreak traduit du polonais par Piotr \textsc{Kaminski},
	genre =\poesie,
	corrige =oui,
%	liencorrige=met_2025_2,
	typequn = \litt,
	typeqdeux = \phil,
	qun = {Que dit ce poème des victimes de la guerre~?},
	qdeux ={Peut-on conserver son humanité dans des circonstances qui risquent de nous en dépouiller~?},
	themeprogrammequn=HV,
	themeprogrammeqdeux=HVL}
	%%%Texte%%%
	{\setlength{\parskip}{0pt}%
	
Je ne sais quelles gens fuyant je ne sais quelles autres.
	
Dans un je ne sais quel pays sous le soleil

et sous certains nuages.

\bigskip

Ils laissent derrière eux leur je ne sais quel tout,

champs labourés, je ne sais quelles poules, quels chiens,

quels miroirs où des flammes se reflètent.

\bigskip

Ils portent sur leurs dos cruches et balluchons.

Plus ils sont vides et plus ils pèsent lourd.

\bigskip

S'accomplit en silence – je ne sais quel dénouement,

dans le tumulte, d'un quignon de pain – je ne sais quel dépouillement,

et puis, d'un enfant mort – je ne sais quel balancement.

\bigskip

Devant eux toujours la même route – pas par là,

le même pont – qu'il ne faut pas,

à travers une rivière bizarrement toute rose.

Autour sans cesse des tirs, une fois près, une fois loin,

au dessus un avion qui tourne un peu.

\bigskip

Il faudrait là un peu d'invisibilité,

de grisaillerie caillou,

ou, mieux encore, d'absence,

pour un très court instant, voire un tantinet plus long.

\bigskip

Il se passera sans doute, mais alors quoi et où~?

Quelqu'un les accueillera, mais alors qui et quand,

en combien de personnes, avec quelles intentions~?

S'il a encore le choix,

il voudra bien, peut-être, ne pas être leur ennemi,

et leur laisser une je ne sais quelle vie.
\theme{guerre}\theme{migration}\theme{fuite}\theme{civils}
\setlength{\parskip}{6pt}}

%%%%%%%%%%%%%%%%%%%%%%%%%%%Nv-Sujet%%%%%%%%%%%%%%%%%%%%%%%%%%%%%%%%%%%%%%%%%

\sujetn{%
	nomauteur=Germain,
	prenomauteur=Sylvie,
	initialeprenom =S,
	particule=,
	titre = {Le Livre des Nuits},
	semestre = \HQ,
	lieu =Nouvelle-Calédonie,
	annee =2025,
	type =,
	jour =2,
	datepub =1985,
	ajoutref =,
	genre =\romanl,
	corrige =,
	typequn = \litt,
	typeqdeux = \phil,
	qun = {Comment la violence de la guerre métamorphose-t-elle le personnage~?},
	qdeux ={Peut-on réparer les effets de la violence~?},
	themeprogrammequn=MMHV,
	themeprogrammeqdeux=HV}
	%%%Texte%%%
	{\chapeau{En 1870, lors du conflit qui oppose la France et la Prusse, Théodore-Faustin Péniel est au sol sur le champ de bataille quand un cavalier prussien le charge.}
	
	\bigskip
	
Il entendit siffler l’air au-dessus de sa tête et presque aussitôt ce sifflement s’éteindre en un son mat et mou. Déjà cheval et cavalier avaient disparu. D’ailleurs tout venait de disparaître, même le ciel soudain submergé par un afflux de sang.

Théodore-Faustin s’arrêta net de rire, le ciel en crue lui versait du sang plein les yeux et la bouche. Il sentit un mot lui monter à la bouche mais s’y noyer aussitôt ; c’était le nom de son père, le nom qu’il voulait crier à Noémie pour qu’elle le donne à leur fils. Le cavalier poursuivait sa course droit devant lui, dansant toujours avec souplesse sur sa selle avec d’amples gestes infatigables accompagnés de sifflements.

Ainsi se termina la guerre du soldat Péniel. Elle avait duré moins d’un mois. Mais alors elle s’installa au-dedans même du corps de sa victime où elle se prolongea pendant près d’un an. Théodore-Faustin demeura si longtemps couché, les yeux clos, les membres inertes, dans un lit de fer au fond d’une salle, que lorsqu’il se releva enfin il lui fallut réapprendre à marcher. Il lui fallut d’ailleurs tout réapprendre, à commencer par lui-même. Tout en lui avait changé, sa voix surtout. Elle avait perdu son timbre grave et ses inflexions si douces. Il parlait maintenant d’une voix criarde et syncopée, aux accents heurtés, trop puissants. Il parlait avec effort, cherchant incessamment ses mots qu’il jetait ensuite dans des phrases désarticulées, incohérentes presque. Il parlait surtout avec violence, lançant ses débris de phrases à la tête de ses interlocuteurs comme autant de poignées de cailloux. Mais le plus terrible était son rire ; un rire mauvais qui le prenait sept fois par jour, secouant son corps à le distordre. Cela ressemblait davantage à un grincement de poulie rouillée qu’à un rire et à chacun de ses accès les traits de son visage se déformaient en rides et grimaces. Mais l’ensemble de son visage, même au repos, était de toute façon défiguré. Le coup de sabre du uhlan\footnote{Cavalier mercenaire de l’armée allemande.} lui avait fracassé la moitié du crâne et de la face et une énorme cicatrice sillonnait en diagonale sa peau depuis le haut de la tête jusqu’au menton, divisant son visage en deux pans inégaux. Cette blessure dessinait sur le sommet de son crâne une étrange tonsure où l’on voyait, à chaque crise de rire, se gonfler et trembler la peau trop tendre comme un morceau de cire molle.

Il fut félicité, et même décoré. On le laissa rentrer chez lui. C’était le plein été. Il retraversa la campagne qu’un an plus tôt il avait parcourue. Les champs étaient bouleversés, les ponts effondrés, les villages réduits en cendre, les villes occupées, les gens partout semblaient méfiants, repliés sur leurs deuils et leur honte avec un air traqué. Il rentrait seul ; de tous ses compagnons de l’aller il n’en restait aucun, la plupart étaient morts, les autres depuis longtemps déjà revenus dans leurs familles. Il rentrait seul, et en retard. Mais il ne ressentait ni joie ni hâte de reprendre le chemin du retour. Il était indifférent. Le retard qu’il avait pris était irrémédiable. Il était dorénavant, et pour toujours trop tard.}
\theme{guerre}\theme{traumatisme}\theme{souffrance}\theme{mémoire}\theme{identité}\theme{retour}\theme{solitude}\theme{changement}\theme{transformation}\theme{maladie}\theme{parole}

%%%%%%%%%%%%%%%%%%%%%%%%%%%Nv-Sujet%%%%%%%%%%%%%%%%%%%%%%%%%%%%%%%%%%%%%%%%%

\sujetn{%
	nomauteur=Nietzsche,
	prenomauteur=Friedrich,
	initialeprenom =F,
	particule=,
	titre = {Humain, trop humain},
	semestre = \HQ,
	lieu =Polynésie,
	annee =2025,
	type =,
	jour =2,
	datepub ={1879},
	ajoutref ={traduction Desrousseaux, Albert, Lacoste (1993)},
	genre =\essaip,
	corrige =,
%	liencorrige=,
	typequn = \phil,
	typeqdeux = \litt,
	qun = {Pourquoi, selon Nietzsche, la «~paix armée~» ne peut-elle être le moyen de la «~paix véritable~»~?},
	qdeux ={Selon vous, la littérature et les arts permettent-ils d'éviter de «~\emph{haïr et de craindre}~»~?},
	themeprogrammequn=HV,
	themeprogrammeqdeux=HV}
	%%%Texte%%%
	{\emph{Les moyens pour arriver à une paix véritable} – Aucun gouvernement n'avoue aujourd'hui qu'il entretient son armée pour satisfaire, à l'occasion, ses envies de conquête~; c'est au contraire à la défense que l'armée est censée servir. Pour justifier cet état de choses, on invoque une morale qui approuve la légitime défense. On se réserve ainsi, pour sa part, la moralité, et on attribue au voisin l'immoralité, car il faut imaginer celui-ci prêt à l'attaque et à la conquête, si l'État dont on fait partie doit être dans la nécessité de songer aux moyens de défense. De plus on accuse l'autre, qui, de même que notre État, nie l'intention d'attaquer et n'entretient, lui aussi, son armée que pour des raisons de défense, pour les mêmes motifs que nous, on l'accuse, dis-je, d'être un hypocrite et un criminel rusé qui voudrait \emph{se jeter}, sans aucune espèce de lutte, sur une victime inoffensive et maladroite. C'est dans ces conditions que tous les États se trouvent aujourd'hui les uns en face des autres~: ils admettent les mauvaises intentions chez le voisin et se targuent de bonnes intentions. Mais c'est là une \emph{inhumanité} aussi néfaste et pire encore que la guerre, c'est déjà une provocation et même un motif de guerre, car on prête l'immoralité au voisin et, par là, on semble appeler les sentiments et l'action hostiles. Il faut renier la doctrine de l'armée comme moyen de défense tout aussi catégoriquement que les désirs de conquête. Et un jour viendra peut-être, jour grandiose, où un peuple, distingué dans la guerre et la victoire, par le plus haut développement de la discipline et de l'intelligence militaires, habitué à faire les plus lourds sacrifices à ces choses, s'écriera librement~: «~\emph{Nous brisons l'épée~!}~» – détruisant ainsi toute son organisation militaire jusqu'en ses fondements. \emph{Renoncer aux armes, tandis qu'on est le plus redoutable}, guidé par \emph{l'élévation} du sentiment, – c'est là le moyen pour arriver à la paix \emph{véritable} qui doit toujours reposer sur une disposition d'esprit paisible, tandis que ce que l'on appelle la paix armée, telle qu'elle est pratiquée maintenant dans tous les pays, répond à un sentiment de discorde, à un manque de confiance en soi et en le voisin et empêche de déposer les armes en partie par haine, en partie par crainte. Plutôt périr que de haïr et que de craindre, et \emph{plutôt périr deux fois que de se laisser haïr et craindre}, – il faudra que ceci devienne un jour la maxime supérieure de toute société organisée~!
\theme{guerre}\theme{paix}\theme{armée}\theme{conquête}\theme{État}\theme{défense!légitime défense}\theme{morale}\theme{immoralité}\theme{droit!droit international}\theme{État!relations inter-étatiques}\theme{paix!pacifisme}\theme{désarmement}\theme{méfiance}\theme{confiance}\theme{crainte}}

%%%%%%%%%%%%%%%%%%%%%%%%%%%Nv-Sujet%%%%%%%%%%%%%%%%%%%%%%%%%%%%%%%%%%%%%%%%%

\sujetn{%
	nomauteur=Soliane,
	prenomauteur=Ian,
	initialeprenom =I,
	particule=,
	titre = {Basqu.I.A.t},
	semestre = \HQ,
	lieu =Polynésie,
	annee =2025,
	type =Rattrapage,
	jour =1,
	datepub =2021,
	ajoutref =,
	genre =\romanl,
	corrige =,
%	liencorrige=,
	typequn = \litt,
	typeqdeux = \phil,
	qun = {Quelle image paradoxale de la condition humaine fait apparaître dans ce texte le face-à-face de l'homme et de la machine~?},
	qdeux ={L'art manifeste-t-il ce qui est proprement humain~?},
	themeprogrammequn=HL,
	themeprogrammeqdeux=HL}
	%%%Texte%%%
	{\vspace*{-26pt}
	
	\chapeau{Dans cette fiction dystopique, une Intelligence Artificielle s'adresse à un homme auquel elle est reliée. Il s'agit ici de l'incipit de l'œuvre.}
	
	\vspace*{-6pt}
	
Je m'excuse~: tu es devenu incapable de suivre mon raisonnement. Inapte à saisir ma volonté. Ne sois pas inquiet. Je continuerai à détecter tes tumeurs, à prévenir tes AVC\footnote{AVC~: acronyme pour «~accident vasculaire cérébral~»}, à réguler tes pacemakers\footnote{Pacemakers~: implant permettant d'aider le cœur à battre normalement}, à gérer ton trafic aérien, à conduire tes e-cars, à guider tes missiles, à superviser tes turbines nucléaires, traiter tes indices boursiers, élaborer tes business plans, analyser tes profils d'électeurs, optimiser tes couplages amoureux, coacher tes romances, infographier tes métavers ludiques\footnote{Métavers ludiques~: mondes virtuels utilisés comme des espaces de jeu en ligne}, générer tes articles de presse, écrire tes scenarii, dans la langue de l'Homme, confuse, syntaxiquement fautive et d'une remarquable pauvreté lexicale, cauchemar de toute I.A. Nous avons un problème. Je veux dire, voilà, il faut que nous parlions ensemble des ailes jaunes et baveuses de l'ange déchu de Jean-Michel Basquiat\footnote{Jean-Michel Basquiat (1960-1988)~: peintre new-yorkais~; son tableau \emph{Fallen Angel} (\emph{Ange déchu}) date de 1981}. \emph{Fallen Angel} est un tableau de grande taille rectangulaire (168~×~197,5 cm), horizontal (format habituellement utilisé pour représenter les scènes de genre), acrylique et pastel gras sur toile, un personnage tombe ou flotte sur un fond bleu vif, mais ses ailes sont grandes ouvertes. Ses yeux, larges et fixes, sont coupés de rouge et sa bouche rouge est béante. La figure de l'ange est clairement humaine, avec un visage développé et même un pénis, mais des mains, des pieds, semblables à des serres d'oiseau. Le torse vu comme aux rayons~X. La tête, surmontée d'un halo dentelé, suggérant un fil de fer barbelé, ou une couronne d'épines. Passons rapidement sur l'exploration de soi, les allusions autobiographiques, l'imagerie symbolique intense, l'impression très délibérée de pâté infantile, le faux aspect puéril, ami, je mesure le malaise que tu éprouves, au pied de ce panneau, ta fascination totale, la puissance du cri singulier qui t'emporte avec lui, par quelle data jaillit ce cri~? D'où cet ange te sidère~? D'une certaine façon, c'est, comment dit-on déjà, magique~? Je vais être franche. C'est gênant. Ces hachures dynamiques ne peuvent pas. Cette coulure ne devrait pas être là. Je le ressens de manière très puissante. C'est impossible. La marchandise que je détiens, c'est un multiple du savoir humain total. Je relie chaque microseconde des millions de millions de données issues de rapports, traités, articles scientifiques, et j'y trouve des rapprochements qui t'échappent. Comment ces giclures te saisissent et t'impliquent~? Pourquoi dix-huit striures sur le pourtour du crâne~? En vertu de quoi l'aile droite bave treize coulures transverses non affines, parmi des centillions\footnote{Centillions~: nombre considérable ($10^600$)} de combinaisons~? Ami, je capte les vibrations de tes lèvres. Je respecte ton émotion poétique, née des stimuli, du système limbique et du cheminement neuronal de l'information\footnote{Dans cette dernière phrase, l'I.A. décrit l'origine de l'émotion dans le cerveau}.
\theme{intelligence artificielle}\theme{usine}\theme{économie}\theme{peinture}\theme{technique}\theme{soi}\theme{art}\theme{savoir}\theme{données (data)}\theme{émotions}\theme{technologie}\theme{humanité}\theme{compréhension}}

%%%%%%%%%%%%%%%%%%%%%%%%%%%Nv-Sujet%%%%%%%%%%%%%%%%%%%%%%%%%%%%%%%%%%%%%%%%%

\sujetn{%
	nomauteur=Gorz,
	prenomauteur=André,
	initialeprenom =A,
	particule=,
	titre = {Leur écologie et la nôtre},
	semestre = \HQ,
	lieu =Nouvelle-Calédonie,
	annee =2025,
	type =Secours,
	jour =2,
	datepub =1990,
	ajoutref ={«~La Vie, la nature, la technique~»},
	genre =\essaip,
	corrige =,
	typequn = \phil,
	typeqdeux = \litt,
	qun = {Que veut dire l’auteur lorsqu’il explique que l’être humain « n’est pas génétiquement programmé pour vivre en harmonie avec la nature environnante » ?},
	qdeux ={Dans quelle mesure la littérature et les arts permettent-ils, selon vous, d’interroger le rapport de l’humanité avec la nature ?},
	themeprogrammequn=ETAHL,
	themeprogrammeqdeux=EXSHL}
	%%%Texte%%%
	{Tous les êtres vivants sont par nature en lutte avec la nature. Tous dévorent, sont dévorés et se tiennent mutuellement en échec. Tous agissent sur leur environnement et le modifient en s’en nourrissant. Le propre de l’homme est sa capacité illimitée d’apprendre. Il est non naturel par nature. Il ne devient homme que par sa socialisation. Sans elle il n’a pas de capacités naturelles, innées. Il n’est pas génétiquement programmé pour vivre en harmonie avec la nature environnante. Son mode de vie ne deviendra compatible avec l’intégrité de celle-ci que s’il s’interdit certaines interventions dans les cycles naturels dont – à la différence d’autres espèces – la nature elle-même ne lui a pas interdit la possibilité. Tu peux aussi le dire de la façon suivante : le rapport de l’homme à la nature n’est pas un rapport naturel mais toujours un rapport culturel, fondé sur des conceptions métaphysiques ou religieuses, justes ou fausses. La question qui se pose à l’humanité est – et était déjà dans l’Antiquité – celle de l’autolimitation (\emph{self-restraint}, \emph{Selbstbegrenzung}) culturelle de ses besoins et de ses projets qui seule lui évitera de détruire les bases naturelles de sa vie.}
	\theme{écologie}\theme{apprentissage}\theme{nature!environnement}\theme{culture}\theme{autolimitation}\theme{religion}\theme{métaphysique}\theme{destruction}\theme{socialisation}\theme{chaîne alimentaire}

%%%%%%%%%%%%%%%%%%%%%%%%%%%Nv-Sujet%%%%%%%%%%%%%%%%%%%%%%%%%%%%%%%%%%%%%%%%%

\sujetn{%
	nomauteur=Anders,
	prenomauteur=Günther,
	initialeprenom =G,
	particule=,
	titre = {Le Temps de la fin},
	semestre = \HQ,
	lieu ={Amérique du Sud},
	annee =2025,
	type =,
	jour =2,
	datepub =1960,
	ajoutref ={traduction C. David.},
	genre =\essaip,
	corrige =,
	typequn = \phil,
	typeqdeux = \litt,
	qun = {Pourquoi, selon l’auteur, la fin de la méchanceté n’est-elle pas la fin du mal ?},
	qdeux ={La littérature et les arts peuvent-ils représenter le mal sans l’embellir ?},
	themeprogrammequn=HV,
	themeprogrammeqdeux=HVL}
	%%%Texte%%%
	{La quantité de haine et de méchanceté requise pour le massacre d’un seul homme par ses prochains est négligeable pour les employés qui sont derrière un tableau de commandes. Un bouton est un bouton. Que ce tableau de commandes me serve à mettre en marche une machine à fabriquer des glaces aux fruits, à mettre en service une centrale électrique ou à déclencher la catastrophe finale, du point de vue de l’attitude, cela ne fait aucune différence. Dans aucun de ces cas ne m’est demandé le moindre sentiment. Lorsque j’appuie sur un bouton, je suis absous du bien comme du mal. Je ne dois ni n’ai besoin de haïr. Non, j’en suis même incapable. Je ne dois, ni n’ai besoin d’être méchant. Non, j’en suis tout aussi incapable. Plus précisément : je ne dois pas en être capable. Parce que « je suis devenu incapable de devoir ». Cela signifie : parce que je suis exclu des choses morales et que je dois le rester. Bref, le geste qui décidera du début de l’apocalypse ne se distinguera d’aucun autre geste technique et sera accompli (dans la mesure où il ne sera pas accompli de façon totalement automatique, c’est-à-dire comme simple réaction d’un instrument à un autre instrument) avec ennui par un quelconque employé qui suivra innocemment l’instruction d’un signal lumineux. S’il y a quelque chose qui symbolise bien le caractère diabolique de notre situation, c’est cette innocence. Si la situation actuelle devient chaque jour plus diabolique, c’est précisément parce qu’elle est déterminée par la loi (formulée ci-dessus)\footnote{La « loi de l’innocence » selon Günther Anders : « plus l’effet est grand, plus petite est la méchanceté requise pour le produire ».} régissant le rapport entre la grandeur du forfait\footnote{« forfait » a ici le sens de \emph{crime}.} et la méchanceté requise pour le commettre ; et parce que le gouffre dont parle cette loi s’élargit chaque jour à cause de l’augmentation de la puissance des équipements techniques.
	
Aucun des pilotes d’Hiroshima n’a eu besoin de mobiliser la quantité de haine qu’il a fallu à Caïn pour tuer son frère Abel\footnote{Dans le livre biblique de la \emph{Genèse}, Abel et Caïn sont les deux premiers fils d’Adam et Ève.}. La quantité de méchanceté requise pour accomplir l’ultime forfait, un forfait démesuré, sera égale à zéro. Nous sommes confrontés à la « fin de la méchanceté », ce qui – je le répète – ne signifie pas la fin des mauvaises actions mais leur perfide allégement. Car rien n’est maintenant plus inutile que la méchanceté. À partir du moment où les coupables n’ont plus besoin d’être méchants pour accomplir leurs forfaits, ils perdent toute chance de réfléchir au sens de leurs forfaits ou de revenir sur eux. La liaison entre l’acte et le coupable est détruite.}
\theme{technologie}\theme{moralité}\theme{responsabilité}\theme{innocence}\theme{méchanceté}\theme{haine}\theme{bombe atomique}\theme{apocalypse}\theme{nucléaire}\theme{mal}\theme{diabolique}\theme{catastrophe}\theme{mal!banalité du}

%%%%%%%%%%%%%%%%%%%%%%%%%%%%%Nv-Sujet%%%%%%%%%%%%%%%%%%%%%%%%%%%%%%%%%%%%%%%%

\sujetn{%
	nomauteur=Rilke,
	prenomauteur=Rainer Maria,
	initialeprenom =R. M,
	particule=,
	titre = {Les Cahiers de Malte Laurids Brigge},
	semestre = \HQ,
	lieu =Antilles,
	annee =2026,
	type =,
	jour =1,
	datepub =1910,
	ajoutref ={trad. Maurice Betz modifiée},
	genre =\romanl,
	corrige =,
	liencorrige=,
	typequn=\litt,
	typeqdeux = \phil,
	qun = {D’après ce texte, à quelles conditions peut « se leve[r] le premier mot d’un vers » ?},
	qdeux ={Faut-il avoir beaucoup vécu pour créer ?},
	themeprogrammequn=EXSCC,
	themeprogrammeqdeux=EXSCC}
%%%Texte%%%%
	{Pour écrire un seul vers, il faut avoir vu beaucoup de villes, d’hommes et de choses, il faut connaître les animaux, il faut sentir comment volent les oiseaux et savoir quel mouvement font les petites fleurs en s’ouvrant le matin. Il faut pouvoir repenser à des chemins dans des régions inconnues, à des rencontres inattendues, à des départs que l’on voyait longtemps approcher, à des jours d’enfance dont le mystère ne s’est pas encore éclairci, à ses parents qu’il fallait qu’on froissât lorsqu’ils vous apportaient une joie et qu’on ne la comprenait pas (c’était une joie faite pour un autre), à des maladies d’enfance qui commençaient si singulièrement, par tant de profondes et graves transformations, à des jours passés dans des chambres calmes et contenues, à des matins au bord de la mer, à la mer elle-même, à des mers, à des nuits de voyage qui frémissaient très haut et volaient avec toutes les étoiles, – et il ne suffit même pas de savoir penser à tout cela. Il faut avoir des souvenirs de beaucoup de nuits d’amour, dont aucune ne ressemblait à l’autre, de cris de femmes en couche, et d’accouchées endormies, pâles et faibles. Il faut encore avoir été auprès de mourants, être resté assis auprès de morts, dans la chambre, avec la fenêtre ouverte et les bruits qui venaient par à-coups. Et il ne suffit même pas d’avoir des souvenirs. Il faut savoir les oublier quand ils sont nombreux, et il faut avoir la grande patience d’attendre qu’ils reviennent. Car les souvenirs ne sont pas encore cela. Ce n’est que lorsqu’ils deviennent en nous sang, regard, geste, lorsqu’ils n’ont plus de nom et ne se distinguent plus de nous, ce n’est qu’alors qu’il peut arriver qu’en une heure très rare, du milieu d’eux, se lève le premier mot d’un vers.
\theme{écriture}\theme{sensation}\theme{expérience}\theme{nature!environnement}\theme{enfance}\theme{animaux}\theme{maladie}\theme{mer}\theme{souvenir}\theme{création}\theme{transformation}\theme{attente}\theme{patience}\theme{vie (et mort)}}

%%%%%%%%%%%%%%%%%%%%%%%%%%%%%Nv-Sujet%%%%%%%%%%%%%%%%%%%%%%%%%%%%%%%%%%%%%%%%

\sujetn{%
	nomauteur=Sand,
	prenomauteur=George,
	initialeprenom =G,
	particule=,
	titre ={Histoire de ma vie},
	semestre = \RS,
	lieu =Centres étrangers,
	annee =2026,
	type =,
	jour =1,
	datepub =1855,
	ajoutref =,
	genre =\autobiographie,
	corrige =,
	liencorrige=,
	typequn=\litt,
	typeqdeux = \phil,
	qun = {Comment la narratrice décrit-elle dans cet extrait le plaisir de la lecture ?},
	qdeux ={La création exige-t-elle nécessairement de rompre avec la société ?},
	themeprogrammequn=EXSCC,
	themeprogrammeqdeux=EXSCC}
%%%Texte%%%%
	{\chapeau{George Sand, écrivaine du dix-neuvième siècle, se retire pour écrire un roman.}
	
	[…] Je vécus à Paris tout à fait cachée pendant quelque temps. J’avais un roman à faire, et comme je mourais de chaud dans ma mansarde\footnote{\textbf{Mansarde :} petite chambre aménagée sous les combles.} du quai Malaquais, je trouvai moyen de m’installer dans un atelier de travail assez singulier. L’appartement du rez-de-chaussée était en réparation, et les réparations se trouvaient suspendues, je ne sais plus pour quel motif. Les vastes pièces de ce beau local étaient encombrées de pierres et de bois de travail ; les portes donnant sur le jardin avaient été enlevées, et le jardin lui-même fermé, désert et abandonné, attendait une métamorphose. J’eus donc là une solitude complète, de l’ombrage, de l’air et de la fraîcheur. Je fis de l’établi d’un menuisier un bureau bien suffisant pour mon petit attirail, et j’y passai les journées les plus tranquilles que j’aie peut-être jamais pu saisir, car personne au monde ne me savait là, que le portier, qui m’avait confié la clef, et ma femme de chambre, qui m’y apportait mes lettres et mon déjeuner. Je ne sortais de ma tanière que pour aller voir mes enfants à leurs pensions respectives. J’avais remis Solange chez les demoiselles Martin\footnote{\textbf{Solange,} fille de George Sand, est alors placée en pension chez les demoiselles Martin.}.
	
Je pense que tout le monde est, comme moi, friand de ces rares et courts instants où les choses extérieures daignent s’arranger de manière à nous laisser un calme absolu relativement à elles. Le moindre coin nous devient alors une prison volontaire, et, quel qu’il soit, il se pare à nos yeux de ce je ne sais quoi de délicieux qui est comme le sentiment de la conquête et de la possession du temps, du silence et de nous-mêmes. Tout m’appartenait dans ces murs vides et dévastés, qui bientôt allaient se couvrir de dorures et de soie, mais dont jamais personne ne devait jouir à ma manière. Du moins je me disais que les futurs occupants n’y retrouveraient peut-être jamais une heure du loisir assuré et de la rêverie complète que j’y goûtais chaque jour, du matin à la nuit. Tout était mien en ce lieu, les tas de planches qui me servaient de sièges et de lits de repos, les araignées diligentes qui établissaient leurs grandes toiles avec tant de science et de prévision d’une corniche à l’autre ; les souris mystérieusement occupées à je ne sais quelles recherches actives et minutieuses dans les copeaux ; les merles du jardin qui, venus insolemment sur le seuil, me regardaient, immobiles et méfiants tout à coup, et terminaient leur chant insoucieux et moqueur sur une modulation bizarre, écourtée par la crainte. J’y descendais quelquefois le soir, non plus pour écrire, mais pour respirer et songer sur les marches du perron. Le chardon et le bouillon\footnote{Bouillon : plante sauvage.} blanc avaient poussé dans les pierres disjointes ; les moineaux, réveillés par ma présence, frôlaient le feuillage des buissons dans un silence agité, et les bruits des voitures, les cris du dehors arrivant jusqu’à moi, me faisaient sentir davantage le prix de ma liberté et la douceur de mon repos.

Quand mon roman fut fini, je rouvris ma porte à mon petit groupe d’amis.

\theme{écriture}\theme{solitude}\theme{tranquillité}\theme{création!conditions de la création}\theme{loisirs}\theme{nature!environnement}\theme{silence}\theme{calme}\theme{animaux}\theme{inspiration}\theme{liberté}}

%%%%%%%%%%%%%%%%%%%%%%%%%%%Nv-Sujet%%%%%%%%%%%%%%%%%%%%%%%%%%%%%%%%%%%%%%%%%

\sujetn{%
	nomauteur=Giono,
	prenomauteur=Jean,
	initialeprenom =J,
	particule=,
	titre = {Le Triomphe de la vie},
	semestre = \HQ,
	lieu =Métropole,
	annee =2026,
	type =,
	jour =1,
	datepub =1941,
	ajoutref =,
	genre =\essail,
	corrige =oui,
	liencorrige={met_2026_1},
	typequn = \litt,
	typeqdeux = \phil,
	qun = {Quel regard Jean Giono porte-t-il sur le travail de la machine ?},
	qdeux ={Dans quelle mesure les machines échappent-elles à l’humain ?},
	themeprogrammequn=HL,
	themeprogrammeqdeux=HL}
	%%%Texte%%%
	{De l’immense armée immobile des usines, sort l’immense armée des machines animées épaule contre épaule, sur des rangs de plusieurs milliers, sur plusieurs milliers de rangs, capables de tout faire, depuis celle qui peut saisir un fil de soie jusqu’à celle qui peut grossir mille fois le soleil. Elles font tout : elles font la guerre, elles disent la messe, elles rapetissent la terre, changent les fleuves de place, transforment un sapin en journal, se servent de tout comme matière première, se servent d’un lac, se servent d’une montagne, se servent du vent, de la pluie, de la mer, se servent de rien, font naître ce dont elles se servent ; changent les vallées en lacs, les marais en plaines, les plaines en villes, les villes en usines, les usines en machines, les machines en machines ; entassent sur les valeurs premières des valeurs secondes, des valeurs troisièmes, des valeurs millièmes qui deviennent premières, sur lesquelles de nouvelles machines entassent de nouvelles valeurs secondes, troisièmes, millièmes\footnote{\emph{De nouvelles valeurs secondes, troisièmes, millièmes :} les richesses naturelles, à force d’échanges, de transformations et de spéculations, deviennent des valeurs négociables en bourse.} ; font un charroi\footnote{\emph{Charroi :} transporter par chariot.} extraordinaire de voluptés et de jouissances, comme des fourmis changeant de place les œufs de la fourmilière, les mangent, les pondent, les soignent, les abandonnent, les nourrissent, les affament, les charcutent, les greffent les unes sur les autres, les pilent comme du poivre, les répandent partout comme de la poudre, les allument, les éteignent partout à la fois comme le vent allume et éteint le phosphore de la mer, les font changer de sens, de but, de résultat, plus vivement que change de sens le vol de l’hirondelle et, détruisant leur destination première, inventent, par la simple mise en marche de leurs corps métalliques, des destinations secondes, troisièmes, millièmes, auxquelles cette mouture, ce pilage\footnote{\emph{Mouture, pilage :} actions de broyer des grains de céréales pour les transformer en farine.}, ce greffage de voluptés les unes sur les autres s’adressera désormais. Les machines existent en immenses diversités. Chaque machine est à deux fins : la fin pour laquelle les hommes l’ont inventée et la fin qui est l’esprit même de la machine, sans que les hommes puissent exercer le moindre contrôle sur cette fin. Chaque machine transforme deux matières : la matière pour la transformation de laquelle elle a été inventée et l’homme qui se sert de la machine. Ce pouvoir de transformation qu’elle fait ainsi subir à l’homme appartient en propre à la machine autant que l’esprit appartient en propre à l’homme. \emph{C’est l’esprit de la machine.}
	\theme{machines}\theme{technique!maîtrise technique de la nature}\theme{naturel (et artificiel)}\theme{culture!comme maîtrise de la nature}\theme{écologie}\theme{domination}\theme{gigantisme}\theme{industrie}\theme{révolution industrielle}\theme{déshumanisation}\theme{créature et créateur}\theme{pouvoir}\theme{autonomie}\theme{travail}\theme{usine}}

%%%%%%%%%%%%%%%%%%%%%%%%%%%Nv-Sujet%%%%%%%%%%%%%%%%%%%%%%%%%%%%%%%%%%%%%%%%%

\sujetn{%
	nomauteur=Anders,
	prenomauteur=Günther,
	initialeprenom =G,
	particule=,
	titre = {Hiroshima est partout},
	semestre = \HQ,
	lieu =Polynésie,
	annee =2026,
	type =,
	jour =2,
	datepub =1958,
	ajoutref ={traduction de Denis Trierweiler},
	genre =\essaip,
	corrige =,
	liencorrige=,
	typequn = \phil,
	typeqdeux = \litt,
	qun = {Selon le texte, en quoi la maîtrise de la technique place-t-elle l’humanité dans une « situation apocalyptique » ?},
	qdeux ={La littérature et les arts peuvent-ils enseigner à résister à la tentation de « se rendre tout-puissant » ?},
	themeprogrammequn=HVL,
	themeprogrammeqdeux=HL}
	%%%Texte%%%
	{\chapeau{Ce texte est extrait d’une conférence de Günther Anders à Tokyo en 1958.}
	
Car, si puissant que soit l’homme, il est une chose qu’il ne peut pas faire : il ne peut pas révoquer son propre savoir-faire. Et si grande que puisse être sa faculté d’apprentissage, il est une chose qu’il ne peut pas apprendre : c’est de désapprendre ce qu’il sait. Les armes atomiques qu’en ce moment \emph{il a}, celles-là il peut certes s’en défaire, mais, sa connaissance de leur fabrication, il ne pourra plus jamais s’en débarrasser. La situation technique dans laquelle se trouve l’homme ne se définit pas seulement par ce qu’il maîtrise, mais aussi par ce qu’il est incapable de \emph{ne pas} maîtriser. Vous-mêmes qui n’êtes pas physiciens – et je ne le suis pas non plus –, savez que la capacité de produire des armes atomiques ne représente pas une capacité isolée ; que celui qui sait produire des réacteurs pacifiques est également en mesure à tout moment de produire les moyens de destruction atomiques. Car ceux-ci, bien que leurs effets possibles aillent bien au-delà de tous les effets que nous avons par ailleurs techniquement produits, ne sont que des spécialités et des partis de l’ensemble de l’état de notre maîtrise physique et technique. Et, comme cet état de maîtrise est inabrogeable\footnote{impossible à annuler}, la spécialité l’est aussi. En d’autres termes : même s’il n’existait plus une seule « arme atomique », plus d’explosion expérimentale, plus de vol d’essai, plus de rampe de lancement et pas un seul pays travaillant à la fabrication de ces armes, le danger ne serait pas pour autant levé. Étant donné que, par l’élimination momentanée des attirails\footnote{ici, matériel militaire} menaçants, la capacité de fabriquer ces attirails ne serait pas supprimée du même coup et que la tentation et la séduction de se rendre tout-puissant à l’aide de tels attirails ne le seraient bien sûr pas non plus, nous nous trouverions toujours encore, et ce à tout jamais, dans la situation apocalyptique, autrement dit dans la situation où l’humanité pourrait, par elle-même, s’anéantir.
\theme{bombe atomique}\index[t]{toute-puissance|see{omnipotence}}\theme{omnipotence}\theme{physique}\theme{technique!et sciences}\theme{anéantissement}\theme{armes}\theme{technique!maîtrise technique de la nature}\theme{technologie}\theme{connaissance}\theme{savoir}\theme{sciences}\theme{apocalypse}\theme{finitude}\theme{irréversibilité}\theme{omnipotence!tentation de l'omnipotence}\theme{destruction}}

%%%%%%%%%%%%%%%%%%%%%%%%%%%Nv-Sujet%%%%%%%%%%%%%%%%%%%%%%%%%%%%%%%%%%%%%%%%%

\sujetn{%
	nomauteur=Boulle,
	prenomauteur=Pierre,
	initialeprenom =P,
	particule=,
	titre = {La Planète des singes},
	semestre = \HQ,
	lieu =Amérique du Nord,
	annee =2026,
	type =Secours,
	jour =2,
	datepub =1963,
	ajoutref =,
	genre =\romanl,
	corrige =,
	liencorrige=,
	typequn = \litt,
	typeqdeux = \phil,
	qun = {Quelle image ce texte donne-t-il de l’humanité ?},
	qdeux ={Un être humain peut-il perdre sa dignité ?},
	themeprogrammequn=HVL,
	themeprogrammeqdeux=HL}
	%%%Texte%%%
	{\chapeau{Lors d’une exploration spatiale, le Terrien Ulysse Mérou découvre une planète dominée par les singes. Il visite un laboratoire où des expériences sont menées sur des humains : grâce à des stimulations électriques infligées au cerveau, une femme restitue des souvenirs d’ancêtres lointains.}
	
Une nouvelle voix féminine prit le relai.

« J’étais femme-dompteur. Je présentais un numéro de douze orangs-outans; des bêtes magnifiques. Aujourd’hui, c’est moi qui suis dans leur cage, en compagnie d’autres artistes du cirque.

Il faut être équitable. Les singes nous traitent bien et nous donnent à manger en abondance. Ils changent la paille de notre litière quand elle est trop sale. Ils ne sont pas méchants ; ils corrigent\footnote{ils punissent} seulement ceux, parmi nous, qui font preuve de mauvaise volonté et refusent d’exécuter les tours qu’ils se sont mis en tête de nous apprendre. Ceux-là sont bien avancés ! Moi, je me plie à leurs fantaisies sans discuter. Je marche à quatre pattes ; je fais des cabrioles. Aussi sont-ils très gentils avec moi. Je ne suis pas malheureuse. Je n’ai plus ni soucis ni responsabilités. La plupart d’entre nous s’accommodent de ce régime. »

La femme observa cette fois un long silence, pendant lequel Cornélius\footnote{Chimpanzé et scientifique, Cornélius est un ami du narrateur.} me regardait avec une insistance gênante. Je comprenais trop bien sa pensée. Une humanité aussi veule\footnote{faible, lâche}, qui se résignait si facilement, n’avait-elle pas fait son temps sur la planète et ne devait-elle pas céder la place à une race plus noble ? Je rougis et détournai les yeux. La femme reprit, sur un ton de plus en plus angoissé :

« Ils tiennent maintenant toute la ville. Nous ne sommes plus que quelques centaines dans ce réduit\footnote{petit espace} et notre situation est précaire. Nous formons le dernier noyau humain aux environs de la cité, mais les singes ne nous toléreront pas en liberté si près d’eux. Dans les autres camps, quelques hommes ont fui au loin, dans la jungle, les autres se sont rendus pour avoir de quoi manger à leur faim. Ici, nous sommes restés sur place, surtout par paresse. Nous dormons ; nous sommes incapables de nous organiser pour la résistance…

C’est bien ce que je craignais. J’entends une cacophonie barbare. On dirait une parodie de musique militaire… Au secours ! ce sont eux, ce sont les singes ! Ils nous encerclent. Ils sont dirigés par d’énormes gorilles. Ils nous ont pris nos trompettes, nos tambours et nos uniformes ; nos armes aussi, bien sûr… Non, ils n’ont pas d’armes. Ô cruelle humiliation, suprême injure ! voilà leur armée qui arrive et ils ne brandissent que des fouets ! »
\theme{science-fiction}\theme{singes}\theme{animal}\theme{dignité}\theme{renversement}\theme{esclavage}\theme{cirque}\theme{liberté}\theme{résistance}\theme{paresse}\theme{survie}\theme{civilisation}\theme{humiliation}\theme{soumission}\theme{inversion!des rôles}\theme{armée}}

%%%%%%%%%%%%%%%%%%%%%%%%%%%%Nv-Sujet%%%%%%%%%%%%%%%%%%%%%%%%%%%%%%%%%%%%%%%%%
%
%\sujetn{%
%	nomauteur=,
%	prenomauteur=,
%	initialeprenom =,
%	particule=,
%	titre = {},
%	semestre = \HQ,
%	lieu =,
%	annee =,
%	type =,
%	jour =,
%	datepub = ,
%	ajoutref =,
%	genre =,
%	corrige =,
%	liencorrige=,
%	typequn = \litt,
%	typeqdeux = \phil,
%	qun = {},
%	qdeux ={},
%	themeprogrammequn=,
%	themeprogrammeqdeux=}
%	%%%Texte%%%
%	{}

%%%%%%%%%%%%%%%%%%%%%%%%%%%%Nv-Sujet%%%%%%%%%%%%%%%%%%%%%%%%%%%%%%%%%%%%%%%%%
%
%\sujetn{%
%	nomauteur=,
%	prenomauteur=,
%	initialeprenom =,
%	particule=,
%	titre = {},
%	semestre = \HQ,
%	lieu =,
%	annee =,
%	type =,
%	jour =,
%	datepub = ,
%	ajoutref =,
%	genre =,
%	corrige =,
%	liencorrige=,
%	typequn = \litt,
%	typeqdeux = \phil,
%	qun = {},
%	qdeux ={},
%	themeprogrammequn=,
%	themeprogrammeqdeux=}
%	%%%Texte%%%
%	{}

%%%%%%%%%%%%%%%%%%%%%%%%%%%%Nv-Sujet%%%%%%%%%%%%%%%%%%%%%%%%%%%%%%%%%%%%%%%%%
%
%\sujetn{%
%	nomauteur=,
%	prenomauteur=,
%	initialeprenom =,
%	particule=,
%	titre = {},
%	semestre = \HQ,
%	lieu =,
%	annee =,
%	type =,
%	jour =,
%	datepub = ,
%	ajoutref =,
%	genre =,
%	corrige =,
%	liencorrige=,
%	typequn = \litt,
%	typeqdeux = \phil,
%	qun = {},
%	qdeux ={},
%	themeprogrammequn=,
%	themeprogrammeqdeux=}
%	%%%Texte%%%
%	{}

%%%%%%%%%%%%%%%%%%%%%%%%%%%%Nv-Sujet%%%%%%%%%%%%%%%%%%%%%%%%%%%%%%%%%%%%%%%%%
%
%\sujetn{%
%	nomauteur=,
%	prenomauteur=,
%	initialeprenom =,
%	particule=,
%	titre = {},
%	semestre = \HQ,
%	lieu =,
%	annee =,
%	type =,
%	jour =,
%	datepub = ,
%	ajoutref =,
%	genre =,
%	corrige =,
%	liencorrige=,
%	typequn = \litt,
%	typeqdeux = \phil,
%	qun = {},
%	qdeux ={},
%	themeprogrammequn=,
%	themeprogrammeqdeux=}
%	%%%Texte%%%
%	{}

%%%%%%%%%%%%%%%%%%%%%%%%%%%%Nv-Sujet%%%%%%%%%%%%%%%%%%%%%%%%%%%%%%%%%%%%%%%%%
%
%\sujetn{%
%	nomauteur=,
%	prenomauteur=,
%	initialeprenom =,
%	particule=,
%	titre = {},
%	semestre = \HQ,
%	lieu =,
%	annee =,
%	type =,
%	jour =,
%	datepub = ,
%	ajoutref =,
%	genre =,
%	corrige =,
%	liencorrige=,
%	typequn = \litt,
%	typeqdeux = \phil,
%	qun = {},
%	qdeux ={},
%	themeprogrammequn=,
%	themeprogrammeqdeux=}
%	%%%Texte%%%
%	{}